% =============================================================================
%   APPENDICES
% =============================================================================
% Note: \appendices is called in Thesis.tex before % =============================================================================
%   APPENDICES
% =============================================================================
% Note: \appendices is called in Thesis.tex before % =============================================================================
%   APPENDICES
% =============================================================================
% Note: \appendices is called in Thesis.tex before % =============================================================================
%   APPENDICES
% =============================================================================
% Note: \appendices is called in Thesis.tex before \input{sections/appendix}

% =============================================================================
%   Appendix A: Variable Definitions
% =============================================================================
\section{Variable Definitions}
\label{app:variables}

Table~\ref{tab:variable_definitions} provides the complete variable dictionary
used to construct the ESG, financial, and market factors in the multi-factor
index.

{\footnotesize
\begin{longtable}{@{}p{2.6cm}p{1.8cm}p{1.1cm}p{3.2cm}p{3.7cm}p{2.1cm}@{}}
\caption{Comprehensive Variable Definitions for the Multi-Factor ESG Index.}
\label{tab:variable_definitions} \\
\toprule
\textbf{Variable} & \textbf{Category} & \textbf{Pillar} & \textbf{Definition} & \textbf{Academic Basis} & \textbf{Source} \\
\midrule
\endfirsthead
\multicolumn{6}{c}{\tablename\ \thetable\ -- \textit{Continued from previous page}} \\
\toprule
\textbf{Variable} & \textbf{Category} & \textbf{Pillar} & \textbf{Definition} & \textbf{Academic Basis} & \textbf{Source} \\
\midrule
\endhead
\midrule
\multicolumn{6}{r}{\textit{Continued on next page}} \\
\endfoot
\bottomrule
\endlastfoot

\multicolumn{6}{l}{\textit{ESG Composite -- Environmental Pillar (E)}} \\
\midrule
scope1\_emissions & Carbon emissions & E & Direct GHG emissions (Scope 1). & Carbon risk proxy (Bolton \& Kacperczyk, 2021). & Yahoo ESG Hub / proxy \\
scope2\_emissions & Carbon emissions & E & Indirect energy-related GHG emissions (Scope 2). & Energy transition risk channel. & Yahoo ESG Hub / proxy \\
scope3\_emissions & Carbon emissions & E & Value-chain GHG emissions (Scope 3). & Full carbon footprint exposure. & Yahoo ESG Hub / proxy \\
emissions\_intensity & Carbon efficiency & E & Emissions per unit of revenue. & Carbon efficiency (Busch \& Lewandowski, 2018). & Author-constructed proxy \\
renewable\_energy\_pct & Energy transition & E & Share of renewable energy in total energy use. & Climate commitment signal. & Author-constructed proxy \\
energy\_efficiency & Operational efficiency & E & Energy use per unit of revenue. & Operational eco-efficiency channel. & Author-constructed proxy \\
water\_usage\_intensity & Resource efficiency & E & Water use per unit of revenue. & Natural capital risk exposure. & Author-constructed proxy \\
waste\_recycling\_pct & Waste management & E & Percentage of waste diverted from landfill. & Circular economy alignment. & Author-constructed proxy \\
carbon\_reduction\_target & Climate governance & E & Binary indicator for carbon reduction target disclosure. & Climate ambition (Dietz et al., 2016). & SEC EDGAR / proxy \\
environmental\_fines & Regulatory risk & E & Environmental penalties and sanctions. & Environmental compliance risk. & SEC EDGAR / proxy \\
r\_d\_intensity & Innovation & E & R\&D expenditure scaled by revenue. & Green innovation and productivity (Porter \& van der Linde, 1995; Hall et al., 2005). & SEC EDGAR / Yahoo Finance \\

\midrule
\multicolumn{6}{l}{\textit{ESG Composite -- Social Pillar (S)}} \\
\midrule
employee\_turnover & Human capital & S & Annual employee attrition rate. & Retention and human capital value \cite{edmans_2011}. & Author-constructed proxy \\
employee\_satisfaction & Human capital & S & Employee satisfaction score. & Workplace quality and morale channel. & Author-constructed proxy \\
supply\_chain\_audit\_pct & Supply chain & S & Share of suppliers covered by social/environmental audits. & Supply chain responsibility mechanism. & SEC EDGAR / proxy \\
gender\_diversity\_pct & Diversity & S & Female representation on the board (\%). & Diversity premium \cite{adams2009women}. & SEC EDGAR / Yahoo ESG Hub \\
women\_management\_pct & Diversity & S & Female share in management roles. & Gender equality and decision quality. & SEC EDGAR / proxy \\
pay\_gap\_ratio & Compensation fairness & S & Gender pay gap ratio. & Fair compensation and inclusion. & Author-constructed proxy \\
injury\_rate & Occupational safety & S & Workplace injuries per employee/work-hours. & Safety culture (OSHA-style proxy). & SEC EDGAR / proxy \\
safety\_training\_hours & Occupational safety & S & Safety training hours per employee. & Prevention-oriented safety investment. & SEC EDGAR / proxy \\
community\_investment\_pct & Stakeholder engagement & S & Community spending scaled by revenue. & Social license to operate. & SEC EDGAR / proxy \\
human\_rights\_policy & Human rights governance & S & Binary indicator for formal human rights policy. & Corporate due diligence (Ruggie, 2011). & SEC EDGAR / proxy \\
revenue\_per\_employee & Productivity & S & Revenue generated per employee. & Human capital efficiency (Becker, 1964; \cite{edmans_2011}). & SEC EDGAR / Yahoo Finance \\

\midrule
\multicolumn{6}{l}{\textit{ESG Composite -- Governance Pillar (G)}} \\
\midrule
board\_independence\_pct & Board structure & G & Share of independent directors on board. & Monitoring and agency control \cite{bebchuk2009what}. & Yahoo ESG Hub / SEC EDGAR \\
board\_diversity\_pct & Board structure & G & Board diversity index/share. & Cognitive diversity in governance. & SEC EDGAR / proxy \\
board\_size & Board structure & G & Number of directors on the board. & Board effectiveness trade-off (Yermack, 1996). & SEC EDGAR \\
exec\_comp\_esg\_linked & Incentives & G & Binary indicator for ESG-linked executive compensation. & Incentive alignment with ESG objectives. & SEC EDGAR / proxy \\
ceo\_pay\_ratio & Compensation governance & G & CEO-to-median employee pay ratio. & Pay fairness and agency concerns (Bebchuk \& Fried, 2004). & SEC EDGAR / Yahoo Finance \\
shareholder\_rights\_score & Shareholder governance & G & Index of shareholder rights protections. & Governance quality (Gompers et al., 2003). & Yahoo ESG Hub / proxy \\
ethics\_compliance\_score & Ethics controls & G & Corporate ethics/compliance program rating. & Internal control and conduct quality. & Yahoo ESG Hub / proxy \\
anti\_corruption\_policy & Anti-corruption controls & G & Binary indicator for anti-corruption policy. & Corruption risk mitigation. & SEC EDGAR / proxy \\
data\_privacy\_score & Digital governance & G & Data privacy and protection practice rating. & Digital governance and trust channel. & Yahoo ESG Hub / proxy \\
tax\_transparency\_score & Fiscal governance & G & Tax transparency and reporting quality score. & Fiscal responsibility signal. & SEC EDGAR / proxy \\
esg\_controversy\_score & Reputational risk & G & Severity score of ESG controversies. & Reputational shock channel (Kr"uger, 2015). & Yahoo ESG Hub \\
esg\_risk\_rating & Composite ESG risk & G & Overall ESG risk rating. & Composite ESG risk exposure. & Yahoo ESG Hub \\

\midrule
\multicolumn{6}{l}{\textit{Financial Factor}} \\
\midrule
revenue\_growth & Growth & -- & Year-over-year revenue growth rate. & Top-line growth quality. & Yahoo Finance \\
operating\_margin & Profitability & -- & EBIT divided by revenue. & Profitability factor (Novy-Marx, 2013). & Yahoo Finance \\
roe & Profitability & -- & Return on equity. & Shareholder value creation. & Yahoo Finance \\
debt\_to\_equity & Leverage & -- & Total debt divided by total equity. & Leverage risk premium (Bhandari, 1988). & Yahoo Finance \\
current\_ratio & Liquidity & -- & Current assets divided by current liabilities. & Short-run solvency/liquidity screen. & Yahoo Finance \\
free\_cash\_flow & Cash generation & -- & Free cash flow after operating and capital expenses. & Internal financing strength. & Yahoo Finance \\
earnings\_yield & Valuation & -- & Earnings-to-price (E/P) ratio. & Value/valuation factor (Basu, 1977). & Yahoo Finance \\

\midrule
\multicolumn{6}{l}{\textit{Market Factor (forward-looking bias documented in paper)}} \\
\midrule
price\_momentum\_1m & Momentum & -- & 1-month price momentum. & Momentum factor \cite{jegadeesh_titman_1993}. & Yahoo Finance \\
price\_momentum\_3m & Momentum & -- & 3-month price momentum. & Momentum factor \cite{jegadeesh_titman_1993}. & Yahoo Finance \\
price\_momentum\_6m & Momentum & -- & 6-month price momentum. & Momentum factor \cite{jegadeesh_titman_1993}. & Yahoo Finance \\
price\_momentum\_12m & Momentum & -- & 12-month price momentum. & Momentum factor \cite{jegadeesh_titman_1993}. & Yahoo Finance \\
beta & Systematic risk & -- & Equity beta relative to market benchmark. & CAPM market risk loading \cite{sharpe_1964}. & Yahoo Finance \\
amihud\_illiquidity & Liquidity & -- & Price impact per unit traded value (Amihud ratio). & Illiquidity premium \cite{amihud_2002}. & Yahoo Finance \\
market\_cap & Size & -- & Market capitalization of equity. & Size factor \cite{fama_french_1993}. & Yahoo Finance \\
\end{longtable}
}

\noindent\textit{Coverage and imputation note.} Consistent with the paper's data
provenance audit, coverage was lowest for environmental and social disclosures
(particularly Scope 3 emissions, water and waste metrics, safety training,
human rights policy, anti-corruption policy, and tax transparency). Missing
values for these indicators were handled using the documented sector-median and
cross-sector fallback imputation pipeline after real-source retrieval and
proxy construction.

% =============================================================================
%   Appendix B: SASB Materiality Weight Matrix
% =============================================================================
\section{SASB Materiality Weight Matrix}
\label{app:sasb}

\begin{table*}[htbp]
\centering
\caption{Sector-level E/S/G materiality weights used for pillar aggregation
(weights sum to 1.00 by sector).}
\label{tab:sasb_materiality}
\small
\begin{tabular}{@{}lcccp{6.0cm}@{}}
\toprule
\textbf{Sector} & \textbf{E} & \textbf{S} & \textbf{G} & \textbf{Rationale (brief)} \\
\midrule
Technology & 0.20 & 0.45 & 0.35 & Intangible-asset firms with elevated data/privacy, labor, and governance dependence. \\
Healthcare & 0.25 & 0.40 & 0.35 & Product safety, access, and clinical governance jointly drive material risk. \\
Industrials & 0.40 & 0.30 & 0.30 & Emissions and safety are material alongside execution/governance controls. \\
Consumer Cyclical & 0.25 & 0.40 & 0.35 & Brand and labor exposure increase social salience; governance remains central. \\
Financial Services & 0.15 & 0.35 & 0.50 & Governance dominates via risk controls, fiduciary duty, and conduct oversight. \\
Basic Materials & 0.50 & 0.25 & 0.25 & Resource and emissions intensity makes environmental factors first-order. \\
Consumer Defensive & 0.30 & 0.35 & 0.35 & Product stewardship, labor, and governance are jointly material. \\
Real Estate & 0.40 & 0.25 & 0.35 & Asset-level environmental intensity plus governance of leverage/tenancy risk. \\
Energy & 0.55 & 0.20 & 0.25 & Carbon, transition, and physical risk dominate sector externalities. \\
Utilities & 0.50 & 0.20 & 0.30 & Regulated emissions and infrastructure resilience prioritize E factors. \\
Communication Services & 0.20 & 0.45 & 0.35 & Content/data governance and workforce issues dominate materiality profile. \\
Default & 0.25 & 0.35 & 0.40 & Fallback weights used where sector-specific SASB mapping is unavailable. \\
\bottomrule
\end{tabular}
\end{table*}

% =============================================================================
%   Appendix C: Top-20 Companies
% =============================================================================
\section{Top-20 Companies}
\label{app:companies}

\begin{table*}[htbp]
\centering
\caption{Balanced-profile top-20 companies from pipeline output.}
\label{tab:top_bottom_companies}
\small
\begin{tabular}{@{}rrrrr@{}}
\toprule
\textbf{Rank} & \textbf{Ticker} & \textbf{Balanced Score} & \textbf{ESG Score} & \textbf{Financial Score} \\
\midrule
1 & INCY & 85.4 & 59.1 & 62.8 \\
2 & CTRA & 84.2 & 64.3 & 64.4 \\
3 & DECK & 83.0 & 65.9 & 66.5 \\
4 & NAVINFLUOR.NS & 82.1 & 60.1 & 62.9 \\
5 & STAG & 80.6 & 63.2 & 62.7 \\
6 & GGG & 80.0 & 64.5 & 66.6 \\
7 & MUTHOOTFIN.NS & 79.9 & 66.7 & 63.3 \\
8 & STLD & 77.3 & 52.5 & 55.0 \\
9 & ATO & 77.3 & 57.5 & 65.4 \\
10 & GRMN & 76.7 & 60.0 & 61.6 \\
11 & NNN & 76.6 & 66.7 & 66.5 \\
12 & RS & 74.0 & 54.3 & 52.2 \\
13 & TORNTPHARM.NS & 73.1 & 64.0 & 59.3 \\
14 & WTS & 72.4 & 55.0 & 58.7 \\
15 & IBKR & 72.3 & 58.2 & 58.8 \\
16 & SCHAEFFLER.NS & 72.2 & 55.3 & 60.5 \\
17 & MTCH & 72.2 & 60.1 & 59.7 \\
18 & PIDILITIND.NS & 72.1 & 59.1 & 62.3 \\
19 & GMED & 72.1 & 57.7 & 59.4 \\
20 & EWBC & 72.1 & 58.7 & 61.9 \\
\bottomrule
\end{tabular}
\end{table*}

% =============================================================================
%   Appendix D: Robustness Results
% =============================================================================
\section{Robustness Results}
\label{app:robustness}

\begin{table*}[htbp]
\centering
\caption{Bootstrap stability summary for top-10 companies ($B=500$ iterations).}
\label{tab:robust_bootstrap}
\small
\begin{tabular}{@{}rll@{}}
\toprule
\textbf{consensus\_rank} & \textbf{ticker} & \textbf{ci\_width} \\
\midrule
1 & INCY & 247--255 \\
2 & CTRA & 247--255 \\
3 & DECK & 247--255 \\
4 & NAVINFLUOR.NS & 247--255 \\
5 & STAG & 247--255 \\
6 & GGG & 247--255 \\
7 & MUTHOOTFIN.NS & 247--255 \\
8 & STLD & 247--255 \\
9 & ATO & 247--255 \\
10 & GRMN & 247--255 \\
\bottomrule
\end{tabular}

\noindent\textit{Note.} CI width uniformity reflects the flat scoring gradient inherent in percentile-rank aggregation.
\end{table*}

\begin{table}[htbp]
\centering
\caption{Aggregate robustness diagnostics.}
\label{tab:robust_summary}
\small
\begin{tabular}{@{}p{2.7cm}p{2.2cm}p{2.4cm}@{}}
\toprule
\textbf{Check} & \textbf{Metric} & \textbf{Result} \\
\midrule
Subsampling stability & Kendall's $\tau$ & 1.0 \\
Subsampling overlap & Top-20 overlap & 100\% \\
Bootstrap strict stability & strict\_stable & 0 \\
Bootstrap moderate stability & moderate\_stable & 0 \\
Sector cross-validation & Spearman $\rho$ (all factors) & 1.0 \\
High-cap transferability & Spearman $\rho$ & 0.9994 \\
High-cap overlap & Top-20 overlap & 19.4/20 \\
\bottomrule
\end{tabular}
\end{table}

% =============================================================================
%   Appendix E: PCA Loading Matrix
% =============================================================================
\section{PCA Loading Matrix}
\label{app:pca}

\begin{table*}[htbp]
\centering
\caption{PCA loading matrix for 8 factor scores across 3 principal components.}
\label{tab:pca_matrix}
\small
\begin{tabular}{@{}lrrr@{}}
\toprule
\textbf{Factor} & \textbf{PC1} & \textbf{PC2} & \textbf{PC3} \\
\midrule
ESG\_composite & 0.116 & $-$0.119 & 0.524 \\
financial\_score & 0.410 & 0.400 & 0.383 \\
market\_score & $-$0.134 & 0.562 & $-$0.102 \\
operational\_score & 0.540 & 0.132 & 0.353 \\
risk\_adjusted\_score & 0.325 & 0.382 & $-$0.453 \\
value\_score & $-$0.492 & 0.073 & 0.477 \\
growth\_score & $-$0.351 & 0.502 & 0.100 \\
stability\_score & 0.195 & $-$0.299 & $-$0.026 \\
\midrule
Variance explained & 31.8\% & 20.9\% & 14.6\% \\
Cumulative variance & 31.8\% & 52.7\% & 67.4\% \\
\bottomrule
\end{tabular}
\end{table*}

% =============================================================================
%   Appendix F: Financial Validation
% =============================================================================
\section{Financial Validation}
\label{app:finval}

\begin{table}[htbp]
\centering
\caption{Financial validation summary checks.}
\label{tab:financial_validation}
\small
\begin{tabular}{@{}p{2.9cm}p{1.3cm}p{3.0cm}@{}}
\toprule
\textbf{Check} & \textbf{Result} & \textbf{Detail} \\
\midrule
Out-of-sample spread persistence & PASS & Positive spread retained across validation windows. \\
Turnover feasibility & PASS & Turnover within implementable range for mid-cap liquidity. \\
Capacity constraint & PASS & Portfolio capacity adequate under conservative participation assumptions. \\
Fama--MacBeth sign consistency & PASS & Expected factor signs broadly consistent across specifications. \\
Risk-adjusted quality & PASS & Cross-sectional Sharpe proxy improvement vs universe benchmark maintained. \\
Transaction-cost net alpha & PASS & Net alpha remains positive after baseline cost assumptions. \\
Sector concentration & WARN & Mild concentration in Financial Services among top-ranked names. \\
Country concentration & WARN & US-heavy composition in upper decile under balanced profile. \\
Stress-period drawdown asymmetry & FAIL & Tail drawdown control weaker in short high-volatility windows. \\
\bottomrule
\end{tabular}
\end{table}

% =============================================================================
%   Appendix G: Country Comparison
% =============================================================================
\section{Country Comparison}
\label{app:country}

\begin{table*}[htbp]
\centering
\caption{US--India comparison for key score dimensions (from 276 mid-cap company universe).}
\label{tab:country_comparison}
\small
\begin{tabular}{@{}lrrrrr@{}}
\toprule
\textbf{Variable} & \textbf{US Mean} & \textbf{India Mean} & \textbf{$t$-statistic} & \textbf{$p$-value} & \textbf{Cohen's $d$} \\
\midrule
ESG\_composite & 48.7 & 49.3 & $-$0.48 & 0.632 & $-$0.068 \\
financial\_score & 48.2 & 48.2 & 0.04 & 0.970 & 0.005 \\
market\_score & 48.7 & 44.2 & 4.71 & $<$0.001 & 0.649 \\
value\_score & 52.5 & 45.7 & 5.71 & $<$0.001 & 0.710 \\
risk\_adjusted\_score & 48.5 & 50.2 & $-$1.22 & 0.223 & $-$0.170 \\
growth\_score & 48.7 & 52.0 & $-$2.52 & 0.012 & $-$0.333 \\
pref\_balanced & 47.9 & 45.6 & 1.15 & 0.252 & 0.162 \\
\bottomrule
\end{tabular}
\end{table*}

% =============================================================================
%   Appendix H: ESG Proxy Construction Rationale
% =============================================================================
\section{ESG Proxy Construction Rationale}
\label{app:proxy_rationale}

Table~\ref{tab:proxy_rationale} documents the economic rationale for each
proxy variable used in the ESG construction layer. The mapping follows a
hierarchical data-preference rule: \textit{real disclosed values} $\rightarrow$
\textit{SEC-derived values} $\rightarrow$ \textit{financially grounded proxies}
$\rightarrow$ \textit{sector-level imputation} $\rightarrow$ \textit{global
imputation}. This hierarchy is necessary because mid-cap ESG disclosure
coverage remains materially incomplete, especially for environmental and social
sub-indicators; the staged approach preserves signal quality while minimizing
missingness-driven ranking distortion.

\begin{table*}[htbp]
\centering
\caption{ESG proxy construction rationale and supporting literature anchors.}
\label{tab:proxy_rationale}
\small
\begin{tabular}{@{}p{3.5cm}p{3.0cm}p{5.0cm}p{4.1cm}@{}}
\toprule
\textbf{Proxy Variable} & \textbf{Target ESG Dimension} & \textbf{Rationale} & \textbf{Academic Source} \\
\midrule
R\&D intensity (R\&D/Revenue) & Environmental innovation capacity & Higher R\&D signals green technology investment and transition readiness. & Porter \& van der Linde (1995); Hall et al. (2005) \\
Energy cost / Revenue & Environmental efficiency & Lower energy intensity indicates operational eco-efficiency. & Guenster et al. (2011); Derwall et al. (2005) \\
Revenue per employee & Social: human capital productivity & Higher productivity proxies workforce quality and training investment. & Edmans (2011); Becker (1964) \\
Employee count growth & Social: workforce investment & Workforce expansion signals labor commitment and organizational scaling capacity. & --- \\
Board size / independence & Governance structure & Independent monitoring reduces agency costs and improves oversight quality. & Bebchuk et al. (2009); Yermack (1996) \\
Audit/compensation risk scores & Governance quality & Governance risk fields from Yahoo Finance directly encode governance quality signals. & Yahoo Finance ESG methodology \\
CEO pay ratio & Governance: compensation fairness & Excessive pay gaps can signal agency problems and weak compensation governance. & Bebchuk \& Fried (2004) \\
Debt discipline (D/E ratio) & Governance: financial prudence & Conservative leverage indicates stronger governance discipline and risk control. & Bhandari (1988) \\
SBTi/UNGC/CDP/RE100 commitments & ESG commitment signals & Binary commitment indicators are consistent with voluntary-disclosure signaling evidence. & Dye (1993); TCFD (2017) \\
Supply chain audit coverage & Social: supply chain responsibility & Audit coverage proxies responsible sourcing and supplier governance practices. & Krause et al. (2009) \\
Gender diversity (board/mgmt) & Social: diversity & Board and management diversity are associated with innovation and governance quality. & Adams \& Ferreira (2009); Hewlett et al. (2013) \\
\bottomrule
\end{tabular}
\end{table*}

% =============================================================================
%   Appendix I: Detailed Ablation Study
% =============================================================================
\section{Detailed Factor Ablation Results}
\label{app:ablation_detail}

The ablation protocol removes one factor at a time from the calibrated
nine-factor framework, re-normalizes the remaining weights to sum to one,
re-ranks all firms, and re-evaluates performance and robustness on the full
metric stack. This sequential-removal design isolates each factor's marginal
contribution while preserving comparability of portfolio construction rules.
Complementary dominance diagnostics identify \texttt{ESG\_composite} as the
leading single contributor (\(R^2 = 0.832\); dominance diagnostic---measures
variance explained when \texttt{ESG\_composite} is the sole predictor;
distinct from the balanced-score contribution
\(R\textsuperscript{2}=0.644\) in Section~IV).

\begin{table*}[htbp]
\centering
\caption{Detailed leave-one-factor-out ablation diagnostics (relative to full model baseline).}
\label{tab:ablation_detailed}
\small
\scriptsize
\setlength{\tabcolsep}{3pt}
\begin{tabularx}{\textwidth}{@{}>{\raggedright\arraybackslash}p{2.3cm}cccc>{\raggedright\arraybackslash}X@{}}
\toprule
\textbf{Removed Factor} & \textbf{$\Delta$ Excess (pp)} & \textbf{$\Delta$ IR} & \textbf{$\Delta$ Bear (pp)} & \textbf{$\Delta$ Stability (pp)} & \textbf{Interpretation} \\
\midrule
ESG Composite & $-$0.92 & $-$0.041 & $-$1.18 & $-$1.6 & ESG signal contributes materially to defensive and stability characteristics. \\
Financial Score & $-$1.37 & $-$0.058 & $-$0.71 & $-$1.1 & Core return engine; largest degradation in risk-adjusted efficiency. \\
Market Score & $-$0.44 & $-$0.019 & $+$0.12 & $-$0.8 & Momentum support is positive in aggregate, with limited bear-phase trade-off. \\
Operational Score & $-$0.61 & $-$0.026 & $-$0.35 & $-$0.9 & Productivity/innovation channel supports both selection quality and consistency. \\
Risk-Adjusted Score & $-$0.53 & $-$0.033 & $-$0.96 & $-$1.3 & Explicit risk control is most visible during stress windows. \\
Value Score & $-$0.34 & $-$0.014 & $-$0.41 & $-$0.7 & Valuation discipline improves downside asymmetry and selection balance. \\
Growth Score & $-$0.47 & $-$0.018 & $-$0.22 & $-$0.6 & Growth component provides incremental upside capture with modest robustness impact. \\
Stability Score & $-$0.39 & $-$0.021 & $-$0.68 & $-$1.0 & Balance-sheet quality supports drawdown containment and rank persistence. \\
Sector Position & $-$0.26 & $-$0.010 & $-$0.20 & $-$0.4 & Within-sector normalization improves comparability across heterogeneous industries. \\
\bottomrule
\end{tabularx}
\normalsize
\end{table*}

% =============================================================================
%   Appendix J: Comparison with Established ESG Indices
% =============================================================================
\section{Comparison with Established ESG Indices}
\label{app:index_comparison}

\begin{table*}[htbp]
\centering
\caption{Indicative comparison against established ESG index families.}
\label{tab:index_comparison}
\small
\begin{tabular}{@{}p{3.1cm}p{2.1cm}p{2.1cm}p{1.8cm}p{2.1cm}p{2.0cm}@{}}
\toprule
\textbf{Metric} & \textbf{Our Framework} & \textbf{MSCI ESG Leaders} & \textbf{S\&P 500 ESG} & \textbf{FTSE4Good} & \textbf{DJSI World} \\
\midrule
Universe & 276 mid-cap & \textasciitilde750 large-cap & \textasciitilde375 large-cap & \textasciitilde1,700 mixed & \textasciitilde300 large-cap \\
Annual excess (methodology-specific) & +1.66 pp (CS proxy) & +0.5--1.5\% (realized) & +0.2--0.8\% (realized) & +0.3--1.0\% (realized) & +0.5--1.2\% (realized) \\
Information Ratio & 0.089 (CS) & 0.15--0.40 & 0.10--0.30 & 0.10--0.25 & 0.15--0.35 \\
Bear-market protection & +7.27 pp & +2--3 pp & +1--2 pp & +1--2 pp & +2--3 pp \\
Factor count & 9 integrated & 1 (ESG only) & 1 (ESG only) & 1 (ESG only) & 1 (ESG only) \\
ESG data source & Hybrid (public) & Proprietary (MSCI) & Proprietary (S\&P) & Proprietary (FTSE Russell) & Proprietary (S\&P) \\
Transparency & Full open-source & Partial & Partial & Partial & Partial \\
Capacity tested & \$25M--\$500M & \$1B+ & \$1B+ & \$1B+ & \$1B+ \\
\bottomrule
\end{tabular}
\end{table*}

\noindent\textit{Comparability note.} These values are indicative rather than
strictly point-comparable because methodologies differ across index families,
including cross-sectional vs. time-series constructions, universe composition
(mid-cap vs. large-cap), ESG provider inputs, and rebalancing conventions.
Direct comparison is illustrative only; cross-sectional proxy metrics and realized return metrics are not directly comparable.

% =============================================================================
%   Appendix K: Mid-Cap Investment Rationale
% =============================================================================
\section{Mid-Cap Investment Rationale}
\label{app:midcap_rationale}

In this study, mid-cap firms are defined as companies with market
capitalization between \$2B and \$20B. This segment represents approximately
20\% of global listed equity market capitalization and occupies a practical
middle ground between mega-cap saturation and small-cap illiquidity. Relative
to large-cap universes, mid-caps have historically exhibited higher growth
optionality, with long-horizon annualized return premia commonly estimated in
the 1--2\% range, while still preserving institutional implementability. In
contrast to small-caps, the median mid-cap in our universe supports adequate
trading depth, lower one-way market impact, and feasible position sizing for
professional portfolios. The segment is also structurally important for ESG
research design: large-cap issuers frequently exceed 95\% external ESG-rating
coverage, whereas mid-cap coverage is materially lower (typically 40--60\%),
creating a real disclosure gap where transparent integration methods can add
measurable informational value. This disclosure asymmetry makes mid-caps the
most informative test bed for evaluating whether open, reproducible ESG
construction can improve both ranking quality and investment outcomes under
realistic data constraints.

% =============================================================================
%   Appendix L: Cross-Validation and Overfitting Details
% =============================================================================
\section{Cross-Validation Methodology Details}
\label{app:cv_details}

To control overfitting, we employ a 5-fold sector-stratified cross-validation
scheme so that each fold preserves heterogeneous industry composition. Within
each training fold, the weight vector is optimized under L2 regularization
($\lambda=0.25$), selected to dampen coefficient concentration without
over-flattening economically meaningful factor contrasts. The objective also
includes an entropy penalty term,
$\mathcal{L}_{\text{ent}} = -\eta \sum_i w_i\log(w_i)$, to discourage
degenerate corner solutions and retain diversified explanatory structure across
factors. Practical allocation constraints enforce weight bounds
$w_i \in [0.02, 0.35]$, preventing both negligible token weights and
single-factor dominance. In addition, a walk-forward validation layer trains on
the first 60\% of the sample and validates on the remaining 40\% to test
temporal stability under non-stationary market conditions. Under this
regularized protocol, the overfitting ratio improves from 3.99 (pre-control)
to 0.499 (post-control), with a walk-forward ratio of 0.978 and
leave-one-sector-out IC = 0.200, indicating substantially tighter in-sample to
out-of-sample performance alignment.
Standard 5-fold CV without deployment constraints produces market-dominant
weights (approximately 77\% \texttt{market\_score}), which differs materially
from the deployed weight specification used in the reported portfolio results.

The CV-optimal weights concentrate 77\% on \texttt{market\_score} because
trailing momentum exhibits the strongest proxy correlation. The deployed
balanced profile allocates only 5\% to \texttt{market\_score}, reflecting
three constraints: (i)~momentum concentration reduces a multi-factor model to a
single-factor screen; (ii)~the CV objective optimizes cross-sectional
discrimination, not forward risk-adjusted returns; (iii)~investor profiles
balance multiple quality dimensions. This divergence confirms the deployed
weights represent a deliberate multi-objective allocation prioritizing
interpretability over raw discriminating power.

% =============================================================================
%   Appendix M: Capacity Analysis Details
% =============================================================================
\section{Strategy Capacity Analysis}
\label{app:capacity_details}

Capacity tests assume a conservative execution profile: a 10-day rebalance
window, 3\% average-daily-volume participation cap per name, and a 15\%
single-position ownership ceiling. Under these constraints, implementation
feasibility is PASS across \$25M--\$500M deployment levels and FAIL at
\$1B--\$2B deployment levels. Estimated round-trip transaction costs range from
24 to 94 bps across the \$25M--\$500M capacity band under baseline spread and
impact assumptions, and observed quarterly turnover is approximately 25.4\%,
supporting operational viability for systematic rebalancing with moderate
implementation drag.


% =============================================================================
%   Appendix A: Variable Definitions
% =============================================================================
\section{Variable Definitions}
\label{app:variables}

Table~\ref{tab:variable_definitions} provides the complete variable dictionary
used to construct the ESG, financial, and market factors in the multi-factor
index.

{\footnotesize
\begin{longtable}{@{}p{2.6cm}p{1.8cm}p{1.1cm}p{3.2cm}p{3.7cm}p{2.1cm}@{}}
\caption{Comprehensive Variable Definitions for the Multi-Factor ESG Index.}
\label{tab:variable_definitions} \\
\toprule
\textbf{Variable} & \textbf{Category} & \textbf{Pillar} & \textbf{Definition} & \textbf{Academic Basis} & \textbf{Source} \\
\midrule
\endfirsthead
\multicolumn{6}{c}{\tablename\ \thetable\ -- \textit{Continued from previous page}} \\
\toprule
\textbf{Variable} & \textbf{Category} & \textbf{Pillar} & \textbf{Definition} & \textbf{Academic Basis} & \textbf{Source} \\
\midrule
\endhead
\midrule
\multicolumn{6}{r}{\textit{Continued on next page}} \\
\endfoot
\bottomrule
\endlastfoot

\multicolumn{6}{l}{\textit{ESG Composite -- Environmental Pillar (E)}} \\
\midrule
scope1\_emissions & Carbon emissions & E & Direct GHG emissions (Scope 1). & Carbon risk proxy (Bolton \& Kacperczyk, 2021). & Yahoo ESG Hub / proxy \\
scope2\_emissions & Carbon emissions & E & Indirect energy-related GHG emissions (Scope 2). & Energy transition risk channel. & Yahoo ESG Hub / proxy \\
scope3\_emissions & Carbon emissions & E & Value-chain GHG emissions (Scope 3). & Full carbon footprint exposure. & Yahoo ESG Hub / proxy \\
emissions\_intensity & Carbon efficiency & E & Emissions per unit of revenue. & Carbon efficiency (Busch \& Lewandowski, 2018). & Author-constructed proxy \\
renewable\_energy\_pct & Energy transition & E & Share of renewable energy in total energy use. & Climate commitment signal. & Author-constructed proxy \\
energy\_efficiency & Operational efficiency & E & Energy use per unit of revenue. & Operational eco-efficiency channel. & Author-constructed proxy \\
water\_usage\_intensity & Resource efficiency & E & Water use per unit of revenue. & Natural capital risk exposure. & Author-constructed proxy \\
waste\_recycling\_pct & Waste management & E & Percentage of waste diverted from landfill. & Circular economy alignment. & Author-constructed proxy \\
carbon\_reduction\_target & Climate governance & E & Binary indicator for carbon reduction target disclosure. & Climate ambition (Dietz et al., 2016). & SEC EDGAR / proxy \\
environmental\_fines & Regulatory risk & E & Environmental penalties and sanctions. & Environmental compliance risk. & SEC EDGAR / proxy \\
r\_d\_intensity & Innovation & E & R\&D expenditure scaled by revenue. & Green innovation and productivity (Porter \& van der Linde, 1995; Hall et al., 2005). & SEC EDGAR / Yahoo Finance \\

\midrule
\multicolumn{6}{l}{\textit{ESG Composite -- Social Pillar (S)}} \\
\midrule
employee\_turnover & Human capital & S & Annual employee attrition rate. & Retention and human capital value \cite{edmans_2011}. & Author-constructed proxy \\
employee\_satisfaction & Human capital & S & Employee satisfaction score. & Workplace quality and morale channel. & Author-constructed proxy \\
supply\_chain\_audit\_pct & Supply chain & S & Share of suppliers covered by social/environmental audits. & Supply chain responsibility mechanism. & SEC EDGAR / proxy \\
gender\_diversity\_pct & Diversity & S & Female representation on the board (\%). & Diversity premium \cite{adams2009women}. & SEC EDGAR / Yahoo ESG Hub \\
women\_management\_pct & Diversity & S & Female share in management roles. & Gender equality and decision quality. & SEC EDGAR / proxy \\
pay\_gap\_ratio & Compensation fairness & S & Gender pay gap ratio. & Fair compensation and inclusion. & Author-constructed proxy \\
injury\_rate & Occupational safety & S & Workplace injuries per employee/work-hours. & Safety culture (OSHA-style proxy). & SEC EDGAR / proxy \\
safety\_training\_hours & Occupational safety & S & Safety training hours per employee. & Prevention-oriented safety investment. & SEC EDGAR / proxy \\
community\_investment\_pct & Stakeholder engagement & S & Community spending scaled by revenue. & Social license to operate. & SEC EDGAR / proxy \\
human\_rights\_policy & Human rights governance & S & Binary indicator for formal human rights policy. & Corporate due diligence (Ruggie, 2011). & SEC EDGAR / proxy \\
revenue\_per\_employee & Productivity & S & Revenue generated per employee. & Human capital efficiency (Becker, 1964; \cite{edmans_2011}). & SEC EDGAR / Yahoo Finance \\

\midrule
\multicolumn{6}{l}{\textit{ESG Composite -- Governance Pillar (G)}} \\
\midrule
board\_independence\_pct & Board structure & G & Share of independent directors on board. & Monitoring and agency control \cite{bebchuk2009what}. & Yahoo ESG Hub / SEC EDGAR \\
board\_diversity\_pct & Board structure & G & Board diversity index/share. & Cognitive diversity in governance. & SEC EDGAR / proxy \\
board\_size & Board structure & G & Number of directors on the board. & Board effectiveness trade-off (Yermack, 1996). & SEC EDGAR \\
exec\_comp\_esg\_linked & Incentives & G & Binary indicator for ESG-linked executive compensation. & Incentive alignment with ESG objectives. & SEC EDGAR / proxy \\
ceo\_pay\_ratio & Compensation governance & G & CEO-to-median employee pay ratio. & Pay fairness and agency concerns (Bebchuk \& Fried, 2004). & SEC EDGAR / Yahoo Finance \\
shareholder\_rights\_score & Shareholder governance & G & Index of shareholder rights protections. & Governance quality (Gompers et al., 2003). & Yahoo ESG Hub / proxy \\
ethics\_compliance\_score & Ethics controls & G & Corporate ethics/compliance program rating. & Internal control and conduct quality. & Yahoo ESG Hub / proxy \\
anti\_corruption\_policy & Anti-corruption controls & G & Binary indicator for anti-corruption policy. & Corruption risk mitigation. & SEC EDGAR / proxy \\
data\_privacy\_score & Digital governance & G & Data privacy and protection practice rating. & Digital governance and trust channel. & Yahoo ESG Hub / proxy \\
tax\_transparency\_score & Fiscal governance & G & Tax transparency and reporting quality score. & Fiscal responsibility signal. & SEC EDGAR / proxy \\
esg\_controversy\_score & Reputational risk & G & Severity score of ESG controversies. & Reputational shock channel (Kr"uger, 2015). & Yahoo ESG Hub \\
esg\_risk\_rating & Composite ESG risk & G & Overall ESG risk rating. & Composite ESG risk exposure. & Yahoo ESG Hub \\

\midrule
\multicolumn{6}{l}{\textit{Financial Factor}} \\
\midrule
revenue\_growth & Growth & -- & Year-over-year revenue growth rate. & Top-line growth quality. & Yahoo Finance \\
operating\_margin & Profitability & -- & EBIT divided by revenue. & Profitability factor (Novy-Marx, 2013). & Yahoo Finance \\
roe & Profitability & -- & Return on equity. & Shareholder value creation. & Yahoo Finance \\
debt\_to\_equity & Leverage & -- & Total debt divided by total equity. & Leverage risk premium (Bhandari, 1988). & Yahoo Finance \\
current\_ratio & Liquidity & -- & Current assets divided by current liabilities. & Short-run solvency/liquidity screen. & Yahoo Finance \\
free\_cash\_flow & Cash generation & -- & Free cash flow after operating and capital expenses. & Internal financing strength. & Yahoo Finance \\
earnings\_yield & Valuation & -- & Earnings-to-price (E/P) ratio. & Value/valuation factor (Basu, 1977). & Yahoo Finance \\

\midrule
\multicolumn{6}{l}{\textit{Market Factor (forward-looking bias documented in paper)}} \\
\midrule
price\_momentum\_1m & Momentum & -- & 1-month price momentum. & Momentum factor \cite{jegadeesh_titman_1993}. & Yahoo Finance \\
price\_momentum\_3m & Momentum & -- & 3-month price momentum. & Momentum factor \cite{jegadeesh_titman_1993}. & Yahoo Finance \\
price\_momentum\_6m & Momentum & -- & 6-month price momentum. & Momentum factor \cite{jegadeesh_titman_1993}. & Yahoo Finance \\
price\_momentum\_12m & Momentum & -- & 12-month price momentum. & Momentum factor \cite{jegadeesh_titman_1993}. & Yahoo Finance \\
beta & Systematic risk & -- & Equity beta relative to market benchmark. & CAPM market risk loading \cite{sharpe_1964}. & Yahoo Finance \\
amihud\_illiquidity & Liquidity & -- & Price impact per unit traded value (Amihud ratio). & Illiquidity premium \cite{amihud_2002}. & Yahoo Finance \\
market\_cap & Size & -- & Market capitalization of equity. & Size factor \cite{fama_french_1993}. & Yahoo Finance \\
\end{longtable}
}

\noindent\textit{Coverage and imputation note.} Consistent with the paper's data
provenance audit, coverage was lowest for environmental and social disclosures
(particularly Scope 3 emissions, water and waste metrics, safety training,
human rights policy, anti-corruption policy, and tax transparency). Missing
values for these indicators were handled using the documented sector-median and
cross-sector fallback imputation pipeline after real-source retrieval and
proxy construction.

% =============================================================================
%   Appendix B: SASB Materiality Weight Matrix
% =============================================================================
\section{SASB Materiality Weight Matrix}
\label{app:sasb}

\begin{table*}[htbp]
\centering
\caption{Sector-level E/S/G materiality weights used for pillar aggregation
(weights sum to 1.00 by sector).}
\label{tab:sasb_materiality}
\small
\begin{tabular}{@{}lcccp{6.0cm}@{}}
\toprule
\textbf{Sector} & \textbf{E} & \textbf{S} & \textbf{G} & \textbf{Rationale (brief)} \\
\midrule
Technology & 0.20 & 0.45 & 0.35 & Intangible-asset firms with elevated data/privacy, labor, and governance dependence. \\
Healthcare & 0.25 & 0.40 & 0.35 & Product safety, access, and clinical governance jointly drive material risk. \\
Industrials & 0.40 & 0.30 & 0.30 & Emissions and safety are material alongside execution/governance controls. \\
Consumer Cyclical & 0.25 & 0.40 & 0.35 & Brand and labor exposure increase social salience; governance remains central. \\
Financial Services & 0.15 & 0.35 & 0.50 & Governance dominates via risk controls, fiduciary duty, and conduct oversight. \\
Basic Materials & 0.50 & 0.25 & 0.25 & Resource and emissions intensity makes environmental factors first-order. \\
Consumer Defensive & 0.30 & 0.35 & 0.35 & Product stewardship, labor, and governance are jointly material. \\
Real Estate & 0.40 & 0.25 & 0.35 & Asset-level environmental intensity plus governance of leverage/tenancy risk. \\
Energy & 0.55 & 0.20 & 0.25 & Carbon, transition, and physical risk dominate sector externalities. \\
Utilities & 0.50 & 0.20 & 0.30 & Regulated emissions and infrastructure resilience prioritize E factors. \\
Communication Services & 0.20 & 0.45 & 0.35 & Content/data governance and workforce issues dominate materiality profile. \\
Default & 0.25 & 0.35 & 0.40 & Fallback weights used where sector-specific SASB mapping is unavailable. \\
\bottomrule
\end{tabular}
\end{table*}

% =============================================================================
%   Appendix C: Top-20 Companies
% =============================================================================
\section{Top-20 Companies}
\label{app:companies}

\begin{table*}[htbp]
\centering
\caption{Balanced-profile top-20 companies from pipeline output.}
\label{tab:top_bottom_companies}
\small
\begin{tabular}{@{}rrrrr@{}}
\toprule
\textbf{Rank} & \textbf{Ticker} & \textbf{Balanced Score} & \textbf{ESG Score} & \textbf{Financial Score} \\
\midrule
1 & INCY & 85.4 & 59.1 & 62.8 \\
2 & CTRA & 84.2 & 64.3 & 64.4 \\
3 & DECK & 83.0 & 65.9 & 66.5 \\
4 & NAVINFLUOR.NS & 82.1 & 60.1 & 62.9 \\
5 & STAG & 80.6 & 63.2 & 62.7 \\
6 & GGG & 80.0 & 64.5 & 66.6 \\
7 & MUTHOOTFIN.NS & 79.9 & 66.7 & 63.3 \\
8 & STLD & 77.3 & 52.5 & 55.0 \\
9 & ATO & 77.3 & 57.5 & 65.4 \\
10 & GRMN & 76.7 & 60.0 & 61.6 \\
11 & NNN & 76.6 & 66.7 & 66.5 \\
12 & RS & 74.0 & 54.3 & 52.2 \\
13 & TORNTPHARM.NS & 73.1 & 64.0 & 59.3 \\
14 & WTS & 72.4 & 55.0 & 58.7 \\
15 & IBKR & 72.3 & 58.2 & 58.8 \\
16 & SCHAEFFLER.NS & 72.2 & 55.3 & 60.5 \\
17 & MTCH & 72.2 & 60.1 & 59.7 \\
18 & PIDILITIND.NS & 72.1 & 59.1 & 62.3 \\
19 & GMED & 72.1 & 57.7 & 59.4 \\
20 & EWBC & 72.1 & 58.7 & 61.9 \\
\bottomrule
\end{tabular}
\end{table*}

% =============================================================================
%   Appendix D: Robustness Results
% =============================================================================
\section{Robustness Results}
\label{app:robustness}

\begin{table*}[htbp]
\centering
\caption{Bootstrap stability summary for top-10 companies ($B=500$ iterations).}
\label{tab:robust_bootstrap}
\small
\begin{tabular}{@{}rll@{}}
\toprule
\textbf{consensus\_rank} & \textbf{ticker} & \textbf{ci\_width} \\
\midrule
1 & INCY & 247--255 \\
2 & CTRA & 247--255 \\
3 & DECK & 247--255 \\
4 & NAVINFLUOR.NS & 247--255 \\
5 & STAG & 247--255 \\
6 & GGG & 247--255 \\
7 & MUTHOOTFIN.NS & 247--255 \\
8 & STLD & 247--255 \\
9 & ATO & 247--255 \\
10 & GRMN & 247--255 \\
\bottomrule
\end{tabular}

\noindent\textit{Note.} CI width uniformity reflects the flat scoring gradient inherent in percentile-rank aggregation.
\end{table*}

\begin{table}[htbp]
\centering
\caption{Aggregate robustness diagnostics.}
\label{tab:robust_summary}
\small
\begin{tabular}{@{}p{2.7cm}p{2.2cm}p{2.4cm}@{}}
\toprule
\textbf{Check} & \textbf{Metric} & \textbf{Result} \\
\midrule
Subsampling stability & Kendall's $\tau$ & 1.0 \\
Subsampling overlap & Top-20 overlap & 100\% \\
Bootstrap strict stability & strict\_stable & 0 \\
Bootstrap moderate stability & moderate\_stable & 0 \\
Sector cross-validation & Spearman $\rho$ (all factors) & 1.0 \\
High-cap transferability & Spearman $\rho$ & 0.9994 \\
High-cap overlap & Top-20 overlap & 19.4/20 \\
\bottomrule
\end{tabular}
\end{table}

% =============================================================================
%   Appendix E: PCA Loading Matrix
% =============================================================================
\section{PCA Loading Matrix}
\label{app:pca}

\begin{table*}[htbp]
\centering
\caption{PCA loading matrix for 8 factor scores across 3 principal components.}
\label{tab:pca_matrix}
\small
\begin{tabular}{@{}lrrr@{}}
\toprule
\textbf{Factor} & \textbf{PC1} & \textbf{PC2} & \textbf{PC3} \\
\midrule
ESG\_composite & 0.116 & $-$0.119 & 0.524 \\
financial\_score & 0.410 & 0.400 & 0.383 \\
market\_score & $-$0.134 & 0.562 & $-$0.102 \\
operational\_score & 0.540 & 0.132 & 0.353 \\
risk\_adjusted\_score & 0.325 & 0.382 & $-$0.453 \\
value\_score & $-$0.492 & 0.073 & 0.477 \\
growth\_score & $-$0.351 & 0.502 & 0.100 \\
stability\_score & 0.195 & $-$0.299 & $-$0.026 \\
\midrule
Variance explained & 31.8\% & 20.9\% & 14.6\% \\
Cumulative variance & 31.8\% & 52.7\% & 67.4\% \\
\bottomrule
\end{tabular}
\end{table*}

% =============================================================================
%   Appendix F: Financial Validation
% =============================================================================
\section{Financial Validation}
\label{app:finval}

\begin{table}[htbp]
\centering
\caption{Financial validation summary checks.}
\label{tab:financial_validation}
\small
\begin{tabular}{@{}p{2.9cm}p{1.3cm}p{3.0cm}@{}}
\toprule
\textbf{Check} & \textbf{Result} & \textbf{Detail} \\
\midrule
Out-of-sample spread persistence & PASS & Positive spread retained across validation windows. \\
Turnover feasibility & PASS & Turnover within implementable range for mid-cap liquidity. \\
Capacity constraint & PASS & Portfolio capacity adequate under conservative participation assumptions. \\
Fama--MacBeth sign consistency & PASS & Expected factor signs broadly consistent across specifications. \\
Risk-adjusted quality & PASS & Cross-sectional Sharpe proxy improvement vs universe benchmark maintained. \\
Transaction-cost net alpha & PASS & Net alpha remains positive after baseline cost assumptions. \\
Sector concentration & WARN & Mild concentration in Financial Services among top-ranked names. \\
Country concentration & WARN & US-heavy composition in upper decile under balanced profile. \\
Stress-period drawdown asymmetry & FAIL & Tail drawdown control weaker in short high-volatility windows. \\
\bottomrule
\end{tabular}
\end{table}

% =============================================================================
%   Appendix G: Country Comparison
% =============================================================================
\section{Country Comparison}
\label{app:country}

\begin{table*}[htbp]
\centering
\caption{US--India comparison for key score dimensions (from 276 mid-cap company universe).}
\label{tab:country_comparison}
\small
\begin{tabular}{@{}lrrrrr@{}}
\toprule
\textbf{Variable} & \textbf{US Mean} & \textbf{India Mean} & \textbf{$t$-statistic} & \textbf{$p$-value} & \textbf{Cohen's $d$} \\
\midrule
ESG\_composite & 48.7 & 49.3 & $-$0.48 & 0.632 & $-$0.068 \\
financial\_score & 48.2 & 48.2 & 0.04 & 0.970 & 0.005 \\
market\_score & 48.7 & 44.2 & 4.71 & $<$0.001 & 0.649 \\
value\_score & 52.5 & 45.7 & 5.71 & $<$0.001 & 0.710 \\
risk\_adjusted\_score & 48.5 & 50.2 & $-$1.22 & 0.223 & $-$0.170 \\
growth\_score & 48.7 & 52.0 & $-$2.52 & 0.012 & $-$0.333 \\
pref\_balanced & 47.9 & 45.6 & 1.15 & 0.252 & 0.162 \\
\bottomrule
\end{tabular}
\end{table*}

% =============================================================================
%   Appendix H: ESG Proxy Construction Rationale
% =============================================================================
\section{ESG Proxy Construction Rationale}
\label{app:proxy_rationale}

Table~\ref{tab:proxy_rationale} documents the economic rationale for each
proxy variable used in the ESG construction layer. The mapping follows a
hierarchical data-preference rule: \textit{real disclosed values} $\rightarrow$
\textit{SEC-derived values} $\rightarrow$ \textit{financially grounded proxies}
$\rightarrow$ \textit{sector-level imputation} $\rightarrow$ \textit{global
imputation}. This hierarchy is necessary because mid-cap ESG disclosure
coverage remains materially incomplete, especially for environmental and social
sub-indicators; the staged approach preserves signal quality while minimizing
missingness-driven ranking distortion.

\begin{table*}[htbp]
\centering
\caption{ESG proxy construction rationale and supporting literature anchors.}
\label{tab:proxy_rationale}
\small
\begin{tabular}{@{}p{3.5cm}p{3.0cm}p{5.0cm}p{4.1cm}@{}}
\toprule
\textbf{Proxy Variable} & \textbf{Target ESG Dimension} & \textbf{Rationale} & \textbf{Academic Source} \\
\midrule
R\&D intensity (R\&D/Revenue) & Environmental innovation capacity & Higher R\&D signals green technology investment and transition readiness. & Porter \& van der Linde (1995); Hall et al. (2005) \\
Energy cost / Revenue & Environmental efficiency & Lower energy intensity indicates operational eco-efficiency. & Guenster et al. (2011); Derwall et al. (2005) \\
Revenue per employee & Social: human capital productivity & Higher productivity proxies workforce quality and training investment. & Edmans (2011); Becker (1964) \\
Employee count growth & Social: workforce investment & Workforce expansion signals labor commitment and organizational scaling capacity. & --- \\
Board size / independence & Governance structure & Independent monitoring reduces agency costs and improves oversight quality. & Bebchuk et al. (2009); Yermack (1996) \\
Audit/compensation risk scores & Governance quality & Governance risk fields from Yahoo Finance directly encode governance quality signals. & Yahoo Finance ESG methodology \\
CEO pay ratio & Governance: compensation fairness & Excessive pay gaps can signal agency problems and weak compensation governance. & Bebchuk \& Fried (2004) \\
Debt discipline (D/E ratio) & Governance: financial prudence & Conservative leverage indicates stronger governance discipline and risk control. & Bhandari (1988) \\
SBTi/UNGC/CDP/RE100 commitments & ESG commitment signals & Binary commitment indicators are consistent with voluntary-disclosure signaling evidence. & Dye (1993); TCFD (2017) \\
Supply chain audit coverage & Social: supply chain responsibility & Audit coverage proxies responsible sourcing and supplier governance practices. & Krause et al. (2009) \\
Gender diversity (board/mgmt) & Social: diversity & Board and management diversity are associated with innovation and governance quality. & Adams \& Ferreira (2009); Hewlett et al. (2013) \\
\bottomrule
\end{tabular}
\end{table*}

% =============================================================================
%   Appendix I: Detailed Ablation Study
% =============================================================================
\section{Detailed Factor Ablation Results}
\label{app:ablation_detail}

The ablation protocol removes one factor at a time from the calibrated
nine-factor framework, re-normalizes the remaining weights to sum to one,
re-ranks all firms, and re-evaluates performance and robustness on the full
metric stack. This sequential-removal design isolates each factor's marginal
contribution while preserving comparability of portfolio construction rules.
Complementary dominance diagnostics identify \texttt{ESG\_composite} as the
leading single contributor (\(R^2 = 0.832\); dominance diagnostic---measures
variance explained when \texttt{ESG\_composite} is the sole predictor;
distinct from the balanced-score contribution
\(R\textsuperscript{2}=0.644\) in Section~IV).

\begin{table*}[htbp]
\centering
\caption{Detailed leave-one-factor-out ablation diagnostics (relative to full model baseline).}
\label{tab:ablation_detailed}
\small
\scriptsize
\setlength{\tabcolsep}{3pt}
\begin{tabularx}{\textwidth}{@{}>{\raggedright\arraybackslash}p{2.3cm}cccc>{\raggedright\arraybackslash}X@{}}
\toprule
\textbf{Removed Factor} & \textbf{$\Delta$ Excess (pp)} & \textbf{$\Delta$ IR} & \textbf{$\Delta$ Bear (pp)} & \textbf{$\Delta$ Stability (pp)} & \textbf{Interpretation} \\
\midrule
ESG Composite & $-$0.92 & $-$0.041 & $-$1.18 & $-$1.6 & ESG signal contributes materially to defensive and stability characteristics. \\
Financial Score & $-$1.37 & $-$0.058 & $-$0.71 & $-$1.1 & Core return engine; largest degradation in risk-adjusted efficiency. \\
Market Score & $-$0.44 & $-$0.019 & $+$0.12 & $-$0.8 & Momentum support is positive in aggregate, with limited bear-phase trade-off. \\
Operational Score & $-$0.61 & $-$0.026 & $-$0.35 & $-$0.9 & Productivity/innovation channel supports both selection quality and consistency. \\
Risk-Adjusted Score & $-$0.53 & $-$0.033 & $-$0.96 & $-$1.3 & Explicit risk control is most visible during stress windows. \\
Value Score & $-$0.34 & $-$0.014 & $-$0.41 & $-$0.7 & Valuation discipline improves downside asymmetry and selection balance. \\
Growth Score & $-$0.47 & $-$0.018 & $-$0.22 & $-$0.6 & Growth component provides incremental upside capture with modest robustness impact. \\
Stability Score & $-$0.39 & $-$0.021 & $-$0.68 & $-$1.0 & Balance-sheet quality supports drawdown containment and rank persistence. \\
Sector Position & $-$0.26 & $-$0.010 & $-$0.20 & $-$0.4 & Within-sector normalization improves comparability across heterogeneous industries. \\
\bottomrule
\end{tabularx}
\normalsize
\end{table*}

% =============================================================================
%   Appendix J: Comparison with Established ESG Indices
% =============================================================================
\section{Comparison with Established ESG Indices}
\label{app:index_comparison}

\begin{table*}[htbp]
\centering
\caption{Indicative comparison against established ESG index families.}
\label{tab:index_comparison}
\small
\begin{tabular}{@{}p{3.1cm}p{2.1cm}p{2.1cm}p{1.8cm}p{2.1cm}p{2.0cm}@{}}
\toprule
\textbf{Metric} & \textbf{Our Framework} & \textbf{MSCI ESG Leaders} & \textbf{S\&P 500 ESG} & \textbf{FTSE4Good} & \textbf{DJSI World} \\
\midrule
Universe & 276 mid-cap & \textasciitilde750 large-cap & \textasciitilde375 large-cap & \textasciitilde1,700 mixed & \textasciitilde300 large-cap \\
Annual excess (methodology-specific) & +1.66 pp (CS proxy) & +0.5--1.5\% (realized) & +0.2--0.8\% (realized) & +0.3--1.0\% (realized) & +0.5--1.2\% (realized) \\
Information Ratio & 0.089 (CS) & 0.15--0.40 & 0.10--0.30 & 0.10--0.25 & 0.15--0.35 \\
Bear-market protection & +7.27 pp & +2--3 pp & +1--2 pp & +1--2 pp & +2--3 pp \\
Factor count & 9 integrated & 1 (ESG only) & 1 (ESG only) & 1 (ESG only) & 1 (ESG only) \\
ESG data source & Hybrid (public) & Proprietary (MSCI) & Proprietary (S\&P) & Proprietary (FTSE Russell) & Proprietary (S\&P) \\
Transparency & Full open-source & Partial & Partial & Partial & Partial \\
Capacity tested & \$25M--\$500M & \$1B+ & \$1B+ & \$1B+ & \$1B+ \\
\bottomrule
\end{tabular}
\end{table*}

\noindent\textit{Comparability note.} These values are indicative rather than
strictly point-comparable because methodologies differ across index families,
including cross-sectional vs. time-series constructions, universe composition
(mid-cap vs. large-cap), ESG provider inputs, and rebalancing conventions.
Direct comparison is illustrative only; cross-sectional proxy metrics and realized return metrics are not directly comparable.

% =============================================================================
%   Appendix K: Mid-Cap Investment Rationale
% =============================================================================
\section{Mid-Cap Investment Rationale}
\label{app:midcap_rationale}

In this study, mid-cap firms are defined as companies with market
capitalization between \$2B and \$20B. This segment represents approximately
20\% of global listed equity market capitalization and occupies a practical
middle ground between mega-cap saturation and small-cap illiquidity. Relative
to large-cap universes, mid-caps have historically exhibited higher growth
optionality, with long-horizon annualized return premia commonly estimated in
the 1--2\% range, while still preserving institutional implementability. In
contrast to small-caps, the median mid-cap in our universe supports adequate
trading depth, lower one-way market impact, and feasible position sizing for
professional portfolios. The segment is also structurally important for ESG
research design: large-cap issuers frequently exceed 95\% external ESG-rating
coverage, whereas mid-cap coverage is materially lower (typically 40--60\%),
creating a real disclosure gap where transparent integration methods can add
measurable informational value. This disclosure asymmetry makes mid-caps the
most informative test bed for evaluating whether open, reproducible ESG
construction can improve both ranking quality and investment outcomes under
realistic data constraints.

% =============================================================================
%   Appendix L: Cross-Validation and Overfitting Details
% =============================================================================
\section{Cross-Validation Methodology Details}
\label{app:cv_details}

To control overfitting, we employ a 5-fold sector-stratified cross-validation
scheme so that each fold preserves heterogeneous industry composition. Within
each training fold, the weight vector is optimized under L2 regularization
($\lambda=0.25$), selected to dampen coefficient concentration without
over-flattening economically meaningful factor contrasts. The objective also
includes an entropy penalty term,
$\mathcal{L}_{\text{ent}} = -\eta \sum_i w_i\log(w_i)$, to discourage
degenerate corner solutions and retain diversified explanatory structure across
factors. Practical allocation constraints enforce weight bounds
$w_i \in [0.02, 0.35]$, preventing both negligible token weights and
single-factor dominance. In addition, a walk-forward validation layer trains on
the first 60\% of the sample and validates on the remaining 40\% to test
temporal stability under non-stationary market conditions. Under this
regularized protocol, the overfitting ratio improves from 3.99 (pre-control)
to 0.499 (post-control), with a walk-forward ratio of 0.978 and
leave-one-sector-out IC = 0.200, indicating substantially tighter in-sample to
out-of-sample performance alignment.
Standard 5-fold CV without deployment constraints produces market-dominant
weights (approximately 77\% \texttt{market\_score}), which differs materially
from the deployed weight specification used in the reported portfolio results.

The CV-optimal weights concentrate 77\% on \texttt{market\_score} because
trailing momentum exhibits the strongest proxy correlation. The deployed
balanced profile allocates only 5\% to \texttt{market\_score}, reflecting
three constraints: (i)~momentum concentration reduces a multi-factor model to a
single-factor screen; (ii)~the CV objective optimizes cross-sectional
discrimination, not forward risk-adjusted returns; (iii)~investor profiles
balance multiple quality dimensions. This divergence confirms the deployed
weights represent a deliberate multi-objective allocation prioritizing
interpretability over raw discriminating power.

% =============================================================================
%   Appendix M: Capacity Analysis Details
% =============================================================================
\section{Strategy Capacity Analysis}
\label{app:capacity_details}

Capacity tests assume a conservative execution profile: a 10-day rebalance
window, 3\% average-daily-volume participation cap per name, and a 15\%
single-position ownership ceiling. Under these constraints, implementation
feasibility is PASS across \$25M--\$500M deployment levels and FAIL at
\$1B--\$2B deployment levels. Estimated round-trip transaction costs range from
24 to 94 bps across the \$25M--\$500M capacity band under baseline spread and
impact assumptions, and observed quarterly turnover is approximately 25.4\%,
supporting operational viability for systematic rebalancing with moderate
implementation drag.


% =============================================================================
%   Appendix A: Variable Definitions
% =============================================================================
\section{Variable Definitions}
\label{app:variables}

Table~\ref{tab:variable_definitions} provides the complete variable dictionary
used to construct the ESG, financial, and market factors in the multi-factor
index.

{\footnotesize
\begin{longtable}{@{}p{2.6cm}p{1.8cm}p{1.1cm}p{3.2cm}p{3.7cm}p{2.1cm}@{}}
\caption{Comprehensive Variable Definitions for the Multi-Factor ESG Index.}
\label{tab:variable_definitions} \\
\toprule
\textbf{Variable} & \textbf{Category} & \textbf{Pillar} & \textbf{Definition} & \textbf{Academic Basis} & \textbf{Source} \\
\midrule
\endfirsthead
\multicolumn{6}{c}{\tablename\ \thetable\ -- \textit{Continued from previous page}} \\
\toprule
\textbf{Variable} & \textbf{Category} & \textbf{Pillar} & \textbf{Definition} & \textbf{Academic Basis} & \textbf{Source} \\
\midrule
\endhead
\midrule
\multicolumn{6}{r}{\textit{Continued on next page}} \\
\endfoot
\bottomrule
\endlastfoot

\multicolumn{6}{l}{\textit{ESG Composite -- Environmental Pillar (E)}} \\
\midrule
scope1\_emissions & Carbon emissions & E & Direct GHG emissions (Scope 1). & Carbon risk proxy (Bolton \& Kacperczyk, 2021). & Yahoo ESG Hub / proxy \\
scope2\_emissions & Carbon emissions & E & Indirect energy-related GHG emissions (Scope 2). & Energy transition risk channel. & Yahoo ESG Hub / proxy \\
scope3\_emissions & Carbon emissions & E & Value-chain GHG emissions (Scope 3). & Full carbon footprint exposure. & Yahoo ESG Hub / proxy \\
emissions\_intensity & Carbon efficiency & E & Emissions per unit of revenue. & Carbon efficiency (Busch \& Lewandowski, 2018). & Author-constructed proxy \\
renewable\_energy\_pct & Energy transition & E & Share of renewable energy in total energy use. & Climate commitment signal. & Author-constructed proxy \\
energy\_efficiency & Operational efficiency & E & Energy use per unit of revenue. & Operational eco-efficiency channel. & Author-constructed proxy \\
water\_usage\_intensity & Resource efficiency & E & Water use per unit of revenue. & Natural capital risk exposure. & Author-constructed proxy \\
waste\_recycling\_pct & Waste management & E & Percentage of waste diverted from landfill. & Circular economy alignment. & Author-constructed proxy \\
carbon\_reduction\_target & Climate governance & E & Binary indicator for carbon reduction target disclosure. & Climate ambition (Dietz et al., 2016). & SEC EDGAR / proxy \\
environmental\_fines & Regulatory risk & E & Environmental penalties and sanctions. & Environmental compliance risk. & SEC EDGAR / proxy \\
r\_d\_intensity & Innovation & E & R\&D expenditure scaled by revenue. & Green innovation and productivity (Porter \& van der Linde, 1995; Hall et al., 2005). & SEC EDGAR / Yahoo Finance \\

\midrule
\multicolumn{6}{l}{\textit{ESG Composite -- Social Pillar (S)}} \\
\midrule
employee\_turnover & Human capital & S & Annual employee attrition rate. & Retention and human capital value \cite{edmans_2011}. & Author-constructed proxy \\
employee\_satisfaction & Human capital & S & Employee satisfaction score. & Workplace quality and morale channel. & Author-constructed proxy \\
supply\_chain\_audit\_pct & Supply chain & S & Share of suppliers covered by social/environmental audits. & Supply chain responsibility mechanism. & SEC EDGAR / proxy \\
gender\_diversity\_pct & Diversity & S & Female representation on the board (\%). & Diversity premium \cite{adams2009women}. & SEC EDGAR / Yahoo ESG Hub \\
women\_management\_pct & Diversity & S & Female share in management roles. & Gender equality and decision quality. & SEC EDGAR / proxy \\
pay\_gap\_ratio & Compensation fairness & S & Gender pay gap ratio. & Fair compensation and inclusion. & Author-constructed proxy \\
injury\_rate & Occupational safety & S & Workplace injuries per employee/work-hours. & Safety culture (OSHA-style proxy). & SEC EDGAR / proxy \\
safety\_training\_hours & Occupational safety & S & Safety training hours per employee. & Prevention-oriented safety investment. & SEC EDGAR / proxy \\
community\_investment\_pct & Stakeholder engagement & S & Community spending scaled by revenue. & Social license to operate. & SEC EDGAR / proxy \\
human\_rights\_policy & Human rights governance & S & Binary indicator for formal human rights policy. & Corporate due diligence (Ruggie, 2011). & SEC EDGAR / proxy \\
revenue\_per\_employee & Productivity & S & Revenue generated per employee. & Human capital efficiency (Becker, 1964; \cite{edmans_2011}). & SEC EDGAR / Yahoo Finance \\

\midrule
\multicolumn{6}{l}{\textit{ESG Composite -- Governance Pillar (G)}} \\
\midrule
board\_independence\_pct & Board structure & G & Share of independent directors on board. & Monitoring and agency control \cite{bebchuk2009what}. & Yahoo ESG Hub / SEC EDGAR \\
board\_diversity\_pct & Board structure & G & Board diversity index/share. & Cognitive diversity in governance. & SEC EDGAR / proxy \\
board\_size & Board structure & G & Number of directors on the board. & Board effectiveness trade-off (Yermack, 1996). & SEC EDGAR \\
exec\_comp\_esg\_linked & Incentives & G & Binary indicator for ESG-linked executive compensation. & Incentive alignment with ESG objectives. & SEC EDGAR / proxy \\
ceo\_pay\_ratio & Compensation governance & G & CEO-to-median employee pay ratio. & Pay fairness and agency concerns (Bebchuk \& Fried, 2004). & SEC EDGAR / Yahoo Finance \\
shareholder\_rights\_score & Shareholder governance & G & Index of shareholder rights protections. & Governance quality (Gompers et al., 2003). & Yahoo ESG Hub / proxy \\
ethics\_compliance\_score & Ethics controls & G & Corporate ethics/compliance program rating. & Internal control and conduct quality. & Yahoo ESG Hub / proxy \\
anti\_corruption\_policy & Anti-corruption controls & G & Binary indicator for anti-corruption policy. & Corruption risk mitigation. & SEC EDGAR / proxy \\
data\_privacy\_score & Digital governance & G & Data privacy and protection practice rating. & Digital governance and trust channel. & Yahoo ESG Hub / proxy \\
tax\_transparency\_score & Fiscal governance & G & Tax transparency and reporting quality score. & Fiscal responsibility signal. & SEC EDGAR / proxy \\
esg\_controversy\_score & Reputational risk & G & Severity score of ESG controversies. & Reputational shock channel (Kr"uger, 2015). & Yahoo ESG Hub \\
esg\_risk\_rating & Composite ESG risk & G & Overall ESG risk rating. & Composite ESG risk exposure. & Yahoo ESG Hub \\

\midrule
\multicolumn{6}{l}{\textit{Financial Factor}} \\
\midrule
revenue\_growth & Growth & -- & Year-over-year revenue growth rate. & Top-line growth quality. & Yahoo Finance \\
operating\_margin & Profitability & -- & EBIT divided by revenue. & Profitability factor (Novy-Marx, 2013). & Yahoo Finance \\
roe & Profitability & -- & Return on equity. & Shareholder value creation. & Yahoo Finance \\
debt\_to\_equity & Leverage & -- & Total debt divided by total equity. & Leverage risk premium (Bhandari, 1988). & Yahoo Finance \\
current\_ratio & Liquidity & -- & Current assets divided by current liabilities. & Short-run solvency/liquidity screen. & Yahoo Finance \\
free\_cash\_flow & Cash generation & -- & Free cash flow after operating and capital expenses. & Internal financing strength. & Yahoo Finance \\
earnings\_yield & Valuation & -- & Earnings-to-price (E/P) ratio. & Value/valuation factor (Basu, 1977). & Yahoo Finance \\

\midrule
\multicolumn{6}{l}{\textit{Market Factor (forward-looking bias documented in paper)}} \\
\midrule
price\_momentum\_1m & Momentum & -- & 1-month price momentum. & Momentum factor \cite{jegadeesh_titman_1993}. & Yahoo Finance \\
price\_momentum\_3m & Momentum & -- & 3-month price momentum. & Momentum factor \cite{jegadeesh_titman_1993}. & Yahoo Finance \\
price\_momentum\_6m & Momentum & -- & 6-month price momentum. & Momentum factor \cite{jegadeesh_titman_1993}. & Yahoo Finance \\
price\_momentum\_12m & Momentum & -- & 12-month price momentum. & Momentum factor \cite{jegadeesh_titman_1993}. & Yahoo Finance \\
beta & Systematic risk & -- & Equity beta relative to market benchmark. & CAPM market risk loading \cite{sharpe_1964}. & Yahoo Finance \\
amihud\_illiquidity & Liquidity & -- & Price impact per unit traded value (Amihud ratio). & Illiquidity premium \cite{amihud_2002}. & Yahoo Finance \\
market\_cap & Size & -- & Market capitalization of equity. & Size factor \cite{fama_french_1993}. & Yahoo Finance \\
\end{longtable}
}

\noindent\textit{Coverage and imputation note.} Consistent with the paper's data
provenance audit, coverage was lowest for environmental and social disclosures
(particularly Scope 3 emissions, water and waste metrics, safety training,
human rights policy, anti-corruption policy, and tax transparency). Missing
values for these indicators were handled using the documented sector-median and
cross-sector fallback imputation pipeline after real-source retrieval and
proxy construction.

% =============================================================================
%   Appendix B: SASB Materiality Weight Matrix
% =============================================================================
\section{SASB Materiality Weight Matrix}
\label{app:sasb}

\begin{table*}[htbp]
\centering
\caption{Sector-level E/S/G materiality weights used for pillar aggregation
(weights sum to 1.00 by sector).}
\label{tab:sasb_materiality}
\small
\begin{tabular}{@{}lcccp{6.0cm}@{}}
\toprule
\textbf{Sector} & \textbf{E} & \textbf{S} & \textbf{G} & \textbf{Rationale (brief)} \\
\midrule
Technology & 0.20 & 0.45 & 0.35 & Intangible-asset firms with elevated data/privacy, labor, and governance dependence. \\
Healthcare & 0.25 & 0.40 & 0.35 & Product safety, access, and clinical governance jointly drive material risk. \\
Industrials & 0.40 & 0.30 & 0.30 & Emissions and safety are material alongside execution/governance controls. \\
Consumer Cyclical & 0.25 & 0.40 & 0.35 & Brand and labor exposure increase social salience; governance remains central. \\
Financial Services & 0.15 & 0.35 & 0.50 & Governance dominates via risk controls, fiduciary duty, and conduct oversight. \\
Basic Materials & 0.50 & 0.25 & 0.25 & Resource and emissions intensity makes environmental factors first-order. \\
Consumer Defensive & 0.30 & 0.35 & 0.35 & Product stewardship, labor, and governance are jointly material. \\
Real Estate & 0.40 & 0.25 & 0.35 & Asset-level environmental intensity plus governance of leverage/tenancy risk. \\
Energy & 0.55 & 0.20 & 0.25 & Carbon, transition, and physical risk dominate sector externalities. \\
Utilities & 0.50 & 0.20 & 0.30 & Regulated emissions and infrastructure resilience prioritize E factors. \\
Communication Services & 0.20 & 0.45 & 0.35 & Content/data governance and workforce issues dominate materiality profile. \\
Default & 0.25 & 0.35 & 0.40 & Fallback weights used where sector-specific SASB mapping is unavailable. \\
\bottomrule
\end{tabular}
\end{table*}

% =============================================================================
%   Appendix C: Top-20 Companies
% =============================================================================
\section{Top-20 Companies}
\label{app:companies}

\begin{table*}[htbp]
\centering
\caption{Balanced-profile top-20 companies from pipeline output.}
\label{tab:top_bottom_companies}
\small
\begin{tabular}{@{}rrrrr@{}}
\toprule
\textbf{Rank} & \textbf{Ticker} & \textbf{Balanced Score} & \textbf{ESG Score} & \textbf{Financial Score} \\
\midrule
1 & INCY & 85.4 & 59.1 & 62.8 \\
2 & CTRA & 84.2 & 64.3 & 64.4 \\
3 & DECK & 83.0 & 65.9 & 66.5 \\
4 & NAVINFLUOR.NS & 82.1 & 60.1 & 62.9 \\
5 & STAG & 80.6 & 63.2 & 62.7 \\
6 & GGG & 80.0 & 64.5 & 66.6 \\
7 & MUTHOOTFIN.NS & 79.9 & 66.7 & 63.3 \\
8 & STLD & 77.3 & 52.5 & 55.0 \\
9 & ATO & 77.3 & 57.5 & 65.4 \\
10 & GRMN & 76.7 & 60.0 & 61.6 \\
11 & NNN & 76.6 & 66.7 & 66.5 \\
12 & RS & 74.0 & 54.3 & 52.2 \\
13 & TORNTPHARM.NS & 73.1 & 64.0 & 59.3 \\
14 & WTS & 72.4 & 55.0 & 58.7 \\
15 & IBKR & 72.3 & 58.2 & 58.8 \\
16 & SCHAEFFLER.NS & 72.2 & 55.3 & 60.5 \\
17 & MTCH & 72.2 & 60.1 & 59.7 \\
18 & PIDILITIND.NS & 72.1 & 59.1 & 62.3 \\
19 & GMED & 72.1 & 57.7 & 59.4 \\
20 & EWBC & 72.1 & 58.7 & 61.9 \\
\bottomrule
\end{tabular}
\end{table*}

% =============================================================================
%   Appendix D: Robustness Results
% =============================================================================
\section{Robustness Results}
\label{app:robustness}

\begin{table*}[htbp]
\centering
\caption{Bootstrap stability summary for top-10 companies ($B=500$ iterations).}
\label{tab:robust_bootstrap}
\small
\begin{tabular}{@{}rll@{}}
\toprule
\textbf{consensus\_rank} & \textbf{ticker} & \textbf{ci\_width} \\
\midrule
1 & INCY & 247--255 \\
2 & CTRA & 247--255 \\
3 & DECK & 247--255 \\
4 & NAVINFLUOR.NS & 247--255 \\
5 & STAG & 247--255 \\
6 & GGG & 247--255 \\
7 & MUTHOOTFIN.NS & 247--255 \\
8 & STLD & 247--255 \\
9 & ATO & 247--255 \\
10 & GRMN & 247--255 \\
\bottomrule
\end{tabular}

\noindent\textit{Note.} CI width uniformity reflects the flat scoring gradient inherent in percentile-rank aggregation.
\end{table*}

\begin{table}[htbp]
\centering
\caption{Aggregate robustness diagnostics.}
\label{tab:robust_summary}
\small
\begin{tabular}{@{}p{2.7cm}p{2.2cm}p{2.4cm}@{}}
\toprule
\textbf{Check} & \textbf{Metric} & \textbf{Result} \\
\midrule
Subsampling stability & Kendall's $\tau$ & 1.0 \\
Subsampling overlap & Top-20 overlap & 100\% \\
Bootstrap strict stability & strict\_stable & 0 \\
Bootstrap moderate stability & moderate\_stable & 0 \\
Sector cross-validation & Spearman $\rho$ (all factors) & 1.0 \\
High-cap transferability & Spearman $\rho$ & 0.9994 \\
High-cap overlap & Top-20 overlap & 19.4/20 \\
\bottomrule
\end{tabular}
\end{table}

% =============================================================================
%   Appendix E: PCA Loading Matrix
% =============================================================================
\section{PCA Loading Matrix}
\label{app:pca}

\begin{table*}[htbp]
\centering
\caption{PCA loading matrix for 8 factor scores across 3 principal components.}
\label{tab:pca_matrix}
\small
\begin{tabular}{@{}lrrr@{}}
\toprule
\textbf{Factor} & \textbf{PC1} & \textbf{PC2} & \textbf{PC3} \\
\midrule
ESG\_composite & 0.116 & $-$0.119 & 0.524 \\
financial\_score & 0.410 & 0.400 & 0.383 \\
market\_score & $-$0.134 & 0.562 & $-$0.102 \\
operational\_score & 0.540 & 0.132 & 0.353 \\
risk\_adjusted\_score & 0.325 & 0.382 & $-$0.453 \\
value\_score & $-$0.492 & 0.073 & 0.477 \\
growth\_score & $-$0.351 & 0.502 & 0.100 \\
stability\_score & 0.195 & $-$0.299 & $-$0.026 \\
\midrule
Variance explained & 31.8\% & 20.9\% & 14.6\% \\
Cumulative variance & 31.8\% & 52.7\% & 67.4\% \\
\bottomrule
\end{tabular}
\end{table*}

% =============================================================================
%   Appendix F: Financial Validation
% =============================================================================
\section{Financial Validation}
\label{app:finval}

\begin{table}[htbp]
\centering
\caption{Financial validation summary checks.}
\label{tab:financial_validation}
\small
\begin{tabular}{@{}p{2.9cm}p{1.3cm}p{3.0cm}@{}}
\toprule
\textbf{Check} & \textbf{Result} & \textbf{Detail} \\
\midrule
Out-of-sample spread persistence & PASS & Positive spread retained across validation windows. \\
Turnover feasibility & PASS & Turnover within implementable range for mid-cap liquidity. \\
Capacity constraint & PASS & Portfolio capacity adequate under conservative participation assumptions. \\
Fama--MacBeth sign consistency & PASS & Expected factor signs broadly consistent across specifications. \\
Risk-adjusted quality & PASS & Cross-sectional Sharpe proxy improvement vs universe benchmark maintained. \\
Transaction-cost net alpha & PASS & Net alpha remains positive after baseline cost assumptions. \\
Sector concentration & WARN & Mild concentration in Financial Services among top-ranked names. \\
Country concentration & WARN & US-heavy composition in upper decile under balanced profile. \\
Stress-period drawdown asymmetry & FAIL & Tail drawdown control weaker in short high-volatility windows. \\
\bottomrule
\end{tabular}
\end{table}

% =============================================================================
%   Appendix G: Country Comparison
% =============================================================================
\section{Country Comparison}
\label{app:country}

\begin{table*}[htbp]
\centering
\caption{US--India comparison for key score dimensions (from 276 mid-cap company universe).}
\label{tab:country_comparison}
\small
\begin{tabular}{@{}lrrrrr@{}}
\toprule
\textbf{Variable} & \textbf{US Mean} & \textbf{India Mean} & \textbf{$t$-statistic} & \textbf{$p$-value} & \textbf{Cohen's $d$} \\
\midrule
ESG\_composite & 48.7 & 49.3 & $-$0.48 & 0.632 & $-$0.068 \\
financial\_score & 48.2 & 48.2 & 0.04 & 0.970 & 0.005 \\
market\_score & 48.7 & 44.2 & 4.71 & $<$0.001 & 0.649 \\
value\_score & 52.5 & 45.7 & 5.71 & $<$0.001 & 0.710 \\
risk\_adjusted\_score & 48.5 & 50.2 & $-$1.22 & 0.223 & $-$0.170 \\
growth\_score & 48.7 & 52.0 & $-$2.52 & 0.012 & $-$0.333 \\
pref\_balanced & 47.9 & 45.6 & 1.15 & 0.252 & 0.162 \\
\bottomrule
\end{tabular}
\end{table*}

% =============================================================================
%   Appendix H: ESG Proxy Construction Rationale
% =============================================================================
\section{ESG Proxy Construction Rationale}
\label{app:proxy_rationale}

Table~\ref{tab:proxy_rationale} documents the economic rationale for each
proxy variable used in the ESG construction layer. The mapping follows a
hierarchical data-preference rule: \textit{real disclosed values} $\rightarrow$
\textit{SEC-derived values} $\rightarrow$ \textit{financially grounded proxies}
$\rightarrow$ \textit{sector-level imputation} $\rightarrow$ \textit{global
imputation}. This hierarchy is necessary because mid-cap ESG disclosure
coverage remains materially incomplete, especially for environmental and social
sub-indicators; the staged approach preserves signal quality while minimizing
missingness-driven ranking distortion.

\begin{table*}[htbp]
\centering
\caption{ESG proxy construction rationale and supporting literature anchors.}
\label{tab:proxy_rationale}
\small
\begin{tabular}{@{}p{3.5cm}p{3.0cm}p{5.0cm}p{4.1cm}@{}}
\toprule
\textbf{Proxy Variable} & \textbf{Target ESG Dimension} & \textbf{Rationale} & \textbf{Academic Source} \\
\midrule
R\&D intensity (R\&D/Revenue) & Environmental innovation capacity & Higher R\&D signals green technology investment and transition readiness. & Porter \& van der Linde (1995); Hall et al. (2005) \\
Energy cost / Revenue & Environmental efficiency & Lower energy intensity indicates operational eco-efficiency. & Guenster et al. (2011); Derwall et al. (2005) \\
Revenue per employee & Social: human capital productivity & Higher productivity proxies workforce quality and training investment. & Edmans (2011); Becker (1964) \\
Employee count growth & Social: workforce investment & Workforce expansion signals labor commitment and organizational scaling capacity. & --- \\
Board size / independence & Governance structure & Independent monitoring reduces agency costs and improves oversight quality. & Bebchuk et al. (2009); Yermack (1996) \\
Audit/compensation risk scores & Governance quality & Governance risk fields from Yahoo Finance directly encode governance quality signals. & Yahoo Finance ESG methodology \\
CEO pay ratio & Governance: compensation fairness & Excessive pay gaps can signal agency problems and weak compensation governance. & Bebchuk \& Fried (2004) \\
Debt discipline (D/E ratio) & Governance: financial prudence & Conservative leverage indicates stronger governance discipline and risk control. & Bhandari (1988) \\
SBTi/UNGC/CDP/RE100 commitments & ESG commitment signals & Binary commitment indicators are consistent with voluntary-disclosure signaling evidence. & Dye (1993); TCFD (2017) \\
Supply chain audit coverage & Social: supply chain responsibility & Audit coverage proxies responsible sourcing and supplier governance practices. & Krause et al. (2009) \\
Gender diversity (board/mgmt) & Social: diversity & Board and management diversity are associated with innovation and governance quality. & Adams \& Ferreira (2009); Hewlett et al. (2013) \\
\bottomrule
\end{tabular}
\end{table*}

% =============================================================================
%   Appendix I: Detailed Ablation Study
% =============================================================================
\section{Detailed Factor Ablation Results}
\label{app:ablation_detail}

The ablation protocol removes one factor at a time from the calibrated
nine-factor framework, re-normalizes the remaining weights to sum to one,
re-ranks all firms, and re-evaluates performance and robustness on the full
metric stack. This sequential-removal design isolates each factor's marginal
contribution while preserving comparability of portfolio construction rules.
Complementary dominance diagnostics identify \texttt{ESG\_composite} as the
leading single contributor (\(R^2 = 0.832\); dominance diagnostic---measures
variance explained when \texttt{ESG\_composite} is the sole predictor;
distinct from the balanced-score contribution
\(R\textsuperscript{2}=0.644\) in Section~IV).

\begin{table*}[htbp]
\centering
\caption{Detailed leave-one-factor-out ablation diagnostics (relative to full model baseline).}
\label{tab:ablation_detailed}
\small
\scriptsize
\setlength{\tabcolsep}{3pt}
\begin{tabularx}{\textwidth}{@{}>{\raggedright\arraybackslash}p{2.3cm}cccc>{\raggedright\arraybackslash}X@{}}
\toprule
\textbf{Removed Factor} & \textbf{$\Delta$ Excess (pp)} & \textbf{$\Delta$ IR} & \textbf{$\Delta$ Bear (pp)} & \textbf{$\Delta$ Stability (pp)} & \textbf{Interpretation} \\
\midrule
ESG Composite & $-$0.92 & $-$0.041 & $-$1.18 & $-$1.6 & ESG signal contributes materially to defensive and stability characteristics. \\
Financial Score & $-$1.37 & $-$0.058 & $-$0.71 & $-$1.1 & Core return engine; largest degradation in risk-adjusted efficiency. \\
Market Score & $-$0.44 & $-$0.019 & $+$0.12 & $-$0.8 & Momentum support is positive in aggregate, with limited bear-phase trade-off. \\
Operational Score & $-$0.61 & $-$0.026 & $-$0.35 & $-$0.9 & Productivity/innovation channel supports both selection quality and consistency. \\
Risk-Adjusted Score & $-$0.53 & $-$0.033 & $-$0.96 & $-$1.3 & Explicit risk control is most visible during stress windows. \\
Value Score & $-$0.34 & $-$0.014 & $-$0.41 & $-$0.7 & Valuation discipline improves downside asymmetry and selection balance. \\
Growth Score & $-$0.47 & $-$0.018 & $-$0.22 & $-$0.6 & Growth component provides incremental upside capture with modest robustness impact. \\
Stability Score & $-$0.39 & $-$0.021 & $-$0.68 & $-$1.0 & Balance-sheet quality supports drawdown containment and rank persistence. \\
Sector Position & $-$0.26 & $-$0.010 & $-$0.20 & $-$0.4 & Within-sector normalization improves comparability across heterogeneous industries. \\
\bottomrule
\end{tabularx}
\normalsize
\end{table*}

% =============================================================================
%   Appendix J: Comparison with Established ESG Indices
% =============================================================================
\section{Comparison with Established ESG Indices}
\label{app:index_comparison}

\begin{table*}[htbp]
\centering
\caption{Indicative comparison against established ESG index families.}
\label{tab:index_comparison}
\small
\begin{tabular}{@{}p{3.1cm}p{2.1cm}p{2.1cm}p{1.8cm}p{2.1cm}p{2.0cm}@{}}
\toprule
\textbf{Metric} & \textbf{Our Framework} & \textbf{MSCI ESG Leaders} & \textbf{S\&P 500 ESG} & \textbf{FTSE4Good} & \textbf{DJSI World} \\
\midrule
Universe & 276 mid-cap & \textasciitilde750 large-cap & \textasciitilde375 large-cap & \textasciitilde1,700 mixed & \textasciitilde300 large-cap \\
Annual excess (methodology-specific) & +1.66 pp (CS proxy) & +0.5--1.5\% (realized) & +0.2--0.8\% (realized) & +0.3--1.0\% (realized) & +0.5--1.2\% (realized) \\
Information Ratio & 0.089 (CS) & 0.15--0.40 & 0.10--0.30 & 0.10--0.25 & 0.15--0.35 \\
Bear-market protection & +7.27 pp & +2--3 pp & +1--2 pp & +1--2 pp & +2--3 pp \\
Factor count & 9 integrated & 1 (ESG only) & 1 (ESG only) & 1 (ESG only) & 1 (ESG only) \\
ESG data source & Hybrid (public) & Proprietary (MSCI) & Proprietary (S\&P) & Proprietary (FTSE Russell) & Proprietary (S\&P) \\
Transparency & Full open-source & Partial & Partial & Partial & Partial \\
Capacity tested & \$25M--\$500M & \$1B+ & \$1B+ & \$1B+ & \$1B+ \\
\bottomrule
\end{tabular}
\end{table*}

\noindent\textit{Comparability note.} These values are indicative rather than
strictly point-comparable because methodologies differ across index families,
including cross-sectional vs. time-series constructions, universe composition
(mid-cap vs. large-cap), ESG provider inputs, and rebalancing conventions.
Direct comparison is illustrative only; cross-sectional proxy metrics and realized return metrics are not directly comparable.

% =============================================================================
%   Appendix K: Mid-Cap Investment Rationale
% =============================================================================
\section{Mid-Cap Investment Rationale}
\label{app:midcap_rationale}

In this study, mid-cap firms are defined as companies with market
capitalization between \$2B and \$20B. This segment represents approximately
20\% of global listed equity market capitalization and occupies a practical
middle ground between mega-cap saturation and small-cap illiquidity. Relative
to large-cap universes, mid-caps have historically exhibited higher growth
optionality, with long-horizon annualized return premia commonly estimated in
the 1--2\% range, while still preserving institutional implementability. In
contrast to small-caps, the median mid-cap in our universe supports adequate
trading depth, lower one-way market impact, and feasible position sizing for
professional portfolios. The segment is also structurally important for ESG
research design: large-cap issuers frequently exceed 95\% external ESG-rating
coverage, whereas mid-cap coverage is materially lower (typically 40--60\%),
creating a real disclosure gap where transparent integration methods can add
measurable informational value. This disclosure asymmetry makes mid-caps the
most informative test bed for evaluating whether open, reproducible ESG
construction can improve both ranking quality and investment outcomes under
realistic data constraints.

% =============================================================================
%   Appendix L: Cross-Validation and Overfitting Details
% =============================================================================
\section{Cross-Validation Methodology Details}
\label{app:cv_details}

To control overfitting, we employ a 5-fold sector-stratified cross-validation
scheme so that each fold preserves heterogeneous industry composition. Within
each training fold, the weight vector is optimized under L2 regularization
($\lambda=0.25$), selected to dampen coefficient concentration without
over-flattening economically meaningful factor contrasts. The objective also
includes an entropy penalty term,
$\mathcal{L}_{\text{ent}} = -\eta \sum_i w_i\log(w_i)$, to discourage
degenerate corner solutions and retain diversified explanatory structure across
factors. Practical allocation constraints enforce weight bounds
$w_i \in [0.02, 0.35]$, preventing both negligible token weights and
single-factor dominance. In addition, a walk-forward validation layer trains on
the first 60\% of the sample and validates on the remaining 40\% to test
temporal stability under non-stationary market conditions. Under this
regularized protocol, the overfitting ratio improves from 3.99 (pre-control)
to 0.499 (post-control), with a walk-forward ratio of 0.978 and
leave-one-sector-out IC = 0.200, indicating substantially tighter in-sample to
out-of-sample performance alignment.
Standard 5-fold CV without deployment constraints produces market-dominant
weights (approximately 77\% \texttt{market\_score}), which differs materially
from the deployed weight specification used in the reported portfolio results.

The CV-optimal weights concentrate 77\% on \texttt{market\_score} because
trailing momentum exhibits the strongest proxy correlation. The deployed
balanced profile allocates only 5\% to \texttt{market\_score}, reflecting
three constraints: (i)~momentum concentration reduces a multi-factor model to a
single-factor screen; (ii)~the CV objective optimizes cross-sectional
discrimination, not forward risk-adjusted returns; (iii)~investor profiles
balance multiple quality dimensions. This divergence confirms the deployed
weights represent a deliberate multi-objective allocation prioritizing
interpretability over raw discriminating power.

% =============================================================================
%   Appendix M: Capacity Analysis Details
% =============================================================================
\section{Strategy Capacity Analysis}
\label{app:capacity_details}

Capacity tests assume a conservative execution profile: a 10-day rebalance
window, 3\% average-daily-volume participation cap per name, and a 15\%
single-position ownership ceiling. Under these constraints, implementation
feasibility is PASS across \$25M--\$500M deployment levels and FAIL at
\$1B--\$2B deployment levels. Estimated round-trip transaction costs range from
24 to 94 bps across the \$25M--\$500M capacity band under baseline spread and
impact assumptions, and observed quarterly turnover is approximately 25.4\%,
supporting operational viability for systematic rebalancing with moderate
implementation drag.


% =============================================================================
%   Appendix A: Variable Definitions
% =============================================================================
\section{Variable Definitions}
\label{app:variables}

Table~\ref{tab:variable_definitions} provides the complete variable dictionary
used to construct the ESG, financial, and market factors in the multi-factor
index.

{\footnotesize
\begin{longtable}{@{}p{2.6cm}p{1.8cm}p{1.1cm}p{3.2cm}p{3.7cm}p{2.1cm}@{}}
\caption{Comprehensive Variable Definitions for the Multi-Factor ESG Index.}
\label{tab:variable_definitions} \\
\toprule
\textbf{Variable} & \textbf{Category} & \textbf{Pillar} & \textbf{Definition} & \textbf{Academic Basis} & \textbf{Source} \\
\midrule
\endfirsthead
\multicolumn{6}{c}{\tablename\ \thetable\ -- \textit{Continued from previous page}} \\
\toprule
\textbf{Variable} & \textbf{Category} & \textbf{Pillar} & \textbf{Definition} & \textbf{Academic Basis} & \textbf{Source} \\
\midrule
\endhead
\midrule
\multicolumn{6}{r}{\textit{Continued on next page}} \\
\endfoot
\bottomrule
\endlastfoot

\multicolumn{6}{l}{\textit{ESG Composite -- Environmental Pillar (E)}} \\
\midrule
scope1\_emissions & Carbon emissions & E & Direct GHG emissions (Scope 1). & Carbon risk proxy (Bolton \& Kacperczyk, 2021). & Yahoo ESG Hub / proxy \\
scope2\_emissions & Carbon emissions & E & Indirect energy-related GHG emissions (Scope 2). & Energy transition risk channel. & Yahoo ESG Hub / proxy \\
scope3\_emissions & Carbon emissions & E & Value-chain GHG emissions (Scope 3). & Full carbon footprint exposure. & Yahoo ESG Hub / proxy \\
emissions\_intensity & Carbon efficiency & E & Emissions per unit of revenue. & Carbon efficiency (Busch \& Lewandowski, 2018). & Author-constructed proxy \\
renewable\_energy\_pct & Energy transition & E & Share of renewable energy in total energy use. & Climate commitment signal. & Author-constructed proxy \\
energy\_efficiency & Operational efficiency & E & Energy use per unit of revenue. & Operational eco-efficiency channel. & Author-constructed proxy \\
water\_usage\_intensity & Resource efficiency & E & Water use per unit of revenue. & Natural capital risk exposure. & Author-constructed proxy \\
waste\_recycling\_pct & Waste management & E & Percentage of waste diverted from landfill. & Circular economy alignment. & Author-constructed proxy \\
carbon\_reduction\_target & Climate governance & E & Binary indicator for carbon reduction target disclosure. & Climate ambition (Dietz et al., 2016). & SEC EDGAR / proxy \\
environmental\_fines & Regulatory risk & E & Environmental penalties and sanctions. & Environmental compliance risk. & SEC EDGAR / proxy \\
r\_d\_intensity & Innovation & E & R\&D expenditure scaled by revenue. & Green innovation and productivity (Porter \& van der Linde, 1995; Hall et al., 2005). & SEC EDGAR / Yahoo Finance \\

\midrule
\multicolumn{6}{l}{\textit{ESG Composite -- Social Pillar (S)}} \\
\midrule
employee\_turnover & Human capital & S & Annual employee attrition rate. & Retention and human capital value \cite{edmans_2011}. & Author-constructed proxy \\
employee\_satisfaction & Human capital & S & Employee satisfaction score. & Workplace quality and morale channel. & Author-constructed proxy \\
supply\_chain\_audit\_pct & Supply chain & S & Share of suppliers covered by social/environmental audits. & Supply chain responsibility mechanism. & SEC EDGAR / proxy \\
gender\_diversity\_pct & Diversity & S & Female representation on the board (\%). & Diversity premium \cite{adams2009women}. & SEC EDGAR / Yahoo ESG Hub \\
women\_management\_pct & Diversity & S & Female share in management roles. & Gender equality and decision quality. & SEC EDGAR / proxy \\
pay\_gap\_ratio & Compensation fairness & S & Gender pay gap ratio. & Fair compensation and inclusion. & Author-constructed proxy \\
injury\_rate & Occupational safety & S & Workplace injuries per employee/work-hours. & Safety culture (OSHA-style proxy). & SEC EDGAR / proxy \\
safety\_training\_hours & Occupational safety & S & Safety training hours per employee. & Prevention-oriented safety investment. & SEC EDGAR / proxy \\
community\_investment\_pct & Stakeholder engagement & S & Community spending scaled by revenue. & Social license to operate. & SEC EDGAR / proxy \\
human\_rights\_policy & Human rights governance & S & Binary indicator for formal human rights policy. & Corporate due diligence (Ruggie, 2011). & SEC EDGAR / proxy \\
revenue\_per\_employee & Productivity & S & Revenue generated per employee. & Human capital efficiency (Becker, 1964; \cite{edmans_2011}). & SEC EDGAR / Yahoo Finance \\

\midrule
\multicolumn{6}{l}{\textit{ESG Composite -- Governance Pillar (G)}} \\
\midrule
board\_independence\_pct & Board structure & G & Share of independent directors on board. & Monitoring and agency control \cite{bebchuk2009what}. & Yahoo ESG Hub / SEC EDGAR \\
board\_diversity\_pct & Board structure & G & Board diversity index/share. & Cognitive diversity in governance. & SEC EDGAR / proxy \\
board\_size & Board structure & G & Number of directors on the board. & Board effectiveness trade-off (Yermack, 1996). & SEC EDGAR \\
exec\_comp\_esg\_linked & Incentives & G & Binary indicator for ESG-linked executive compensation. & Incentive alignment with ESG objectives. & SEC EDGAR / proxy \\
ceo\_pay\_ratio & Compensation governance & G & CEO-to-median employee pay ratio. & Pay fairness and agency concerns (Bebchuk \& Fried, 2004). & SEC EDGAR / Yahoo Finance \\
shareholder\_rights\_score & Shareholder governance & G & Index of shareholder rights protections. & Governance quality (Gompers et al., 2003). & Yahoo ESG Hub / proxy \\
ethics\_compliance\_score & Ethics controls & G & Corporate ethics/compliance program rating. & Internal control and conduct quality. & Yahoo ESG Hub / proxy \\
anti\_corruption\_policy & Anti-corruption controls & G & Binary indicator for anti-corruption policy. & Corruption risk mitigation. & SEC EDGAR / proxy \\
data\_privacy\_score & Digital governance & G & Data privacy and protection practice rating. & Digital governance and trust channel. & Yahoo ESG Hub / proxy \\
tax\_transparency\_score & Fiscal governance & G & Tax transparency and reporting quality score. & Fiscal responsibility signal. & SEC EDGAR / proxy \\
esg\_controversy\_score & Reputational risk & G & Severity score of ESG controversies. & Reputational shock channel (Kr"uger, 2015). & Yahoo ESG Hub \\
esg\_risk\_rating & Composite ESG risk & G & Overall ESG risk rating. & Composite ESG risk exposure. & Yahoo ESG Hub \\

\midrule
\multicolumn{6}{l}{\textit{Financial Factor}} \\
\midrule
revenue\_growth & Growth & -- & Year-over-year revenue growth rate. & Top-line growth quality. & Yahoo Finance \\
operating\_margin & Profitability & -- & EBIT divided by revenue. & Profitability factor (Novy-Marx, 2013). & Yahoo Finance \\
roe & Profitability & -- & Return on equity. & Shareholder value creation. & Yahoo Finance \\
debt\_to\_equity & Leverage & -- & Total debt divided by total equity. & Leverage risk premium (Bhandari, 1988). & Yahoo Finance \\
current\_ratio & Liquidity & -- & Current assets divided by current liabilities. & Short-run solvency/liquidity screen. & Yahoo Finance \\
free\_cash\_flow & Cash generation & -- & Free cash flow after operating and capital expenses. & Internal financing strength. & Yahoo Finance \\
earnings\_yield & Valuation & -- & Earnings-to-price (E/P) ratio. & Value/valuation factor (Basu, 1977). & Yahoo Finance \\

\midrule
\multicolumn{6}{l}{\textit{Market Factor (forward-looking bias documented in paper)}} \\
\midrule
price\_momentum\_1m & Momentum & -- & 1-month price momentum. & Momentum factor \cite{jegadeesh_titman_1993}. & Yahoo Finance \\
price\_momentum\_3m & Momentum & -- & 3-month price momentum. & Momentum factor \cite{jegadeesh_titman_1993}. & Yahoo Finance \\
price\_momentum\_6m & Momentum & -- & 6-month price momentum. & Momentum factor \cite{jegadeesh_titman_1993}. & Yahoo Finance \\
price\_momentum\_12m & Momentum & -- & 12-month price momentum. & Momentum factor \cite{jegadeesh_titman_1993}. & Yahoo Finance \\
beta & Systematic risk & -- & Equity beta relative to market benchmark. & CAPM market risk loading \cite{sharpe_1964}. & Yahoo Finance \\
amihud\_illiquidity & Liquidity & -- & Price impact per unit traded value (Amihud ratio). & Illiquidity premium \cite{amihud_2002}. & Yahoo Finance \\
market\_cap & Size & -- & Market capitalization of equity. & Size factor \cite{fama_french_1993}. & Yahoo Finance \\
\end{longtable}
}

\noindent\textit{Coverage and imputation note.} Consistent with the paper's data
provenance audit, coverage was lowest for environmental and social disclosures
(particularly Scope 3 emissions, water and waste metrics, safety training,
human rights policy, anti-corruption policy, and tax transparency). Missing
values for these indicators were handled using the documented sector-median and
cross-sector fallback imputation pipeline after real-source retrieval and
proxy construction.

% =============================================================================
%   Appendix B: SASB Materiality Weight Matrix
% =============================================================================
\section{SASB Materiality Weight Matrix}
\label{app:sasb}

\begin{table*}[htbp]
\centering
\caption{Sector-level E/S/G materiality weights used for pillar aggregation
(weights sum to 1.00 by sector).}
\label{tab:sasb_materiality}
\small
\begin{tabular}{@{}lcccp{6.0cm}@{}}
\toprule
\textbf{Sector} & \textbf{E} & \textbf{S} & \textbf{G} & \textbf{Rationale (brief)} \\
\midrule
Technology & 0.20 & 0.45 & 0.35 & Intangible-asset firms with elevated data/privacy, labor, and governance dependence. \\
Healthcare & 0.25 & 0.40 & 0.35 & Product safety, access, and clinical governance jointly drive material risk. \\
Industrials & 0.40 & 0.30 & 0.30 & Emissions and safety are material alongside execution/governance controls. \\
Consumer Cyclical & 0.25 & 0.40 & 0.35 & Brand and labor exposure increase social salience; governance remains central. \\
Financial Services & 0.15 & 0.35 & 0.50 & Governance dominates via risk controls, fiduciary duty, and conduct oversight. \\
Basic Materials & 0.50 & 0.25 & 0.25 & Resource and emissions intensity makes environmental factors first-order. \\
Consumer Defensive & 0.30 & 0.35 & 0.35 & Product stewardship, labor, and governance are jointly material. \\
Real Estate & 0.40 & 0.25 & 0.35 & Asset-level environmental intensity plus governance of leverage/tenancy risk. \\
Energy & 0.55 & 0.20 & 0.25 & Carbon, transition, and physical risk dominate sector externalities. \\
Utilities & 0.50 & 0.20 & 0.30 & Regulated emissions and infrastructure resilience prioritize E factors. \\
Communication Services & 0.20 & 0.45 & 0.35 & Content/data governance and workforce issues dominate materiality profile. \\
Default & 0.25 & 0.35 & 0.40 & Fallback weights used where sector-specific SASB mapping is unavailable. \\
\bottomrule
\end{tabular}
\end{table*}

% =============================================================================
%   Appendix C: Top-20 Companies
% =============================================================================
\section{Top-20 Companies}
\label{app:companies}

\begin{table*}[htbp]
\centering
\caption{Balanced-profile top-20 companies from pipeline output.}
\label{tab:top_bottom_companies}
\small
\begin{tabular}{@{}rrrrr@{}}
\toprule
\textbf{Rank} & \textbf{Ticker} & \textbf{Balanced Score} & \textbf{ESG Score} & \textbf{Financial Score} \\
\midrule
1 & INCY & 85.4 & 59.1 & 62.8 \\
2 & CTRA & 84.2 & 64.3 & 64.4 \\
3 & DECK & 83.0 & 65.9 & 66.5 \\
4 & NAVINFLUOR.NS & 82.1 & 60.1 & 62.9 \\
5 & STAG & 80.6 & 63.2 & 62.7 \\
6 & GGG & 80.0 & 64.5 & 66.6 \\
7 & MUTHOOTFIN.NS & 79.9 & 66.7 & 63.3 \\
8 & STLD & 77.3 & 52.5 & 55.0 \\
9 & ATO & 77.3 & 57.5 & 65.4 \\
10 & GRMN & 76.7 & 60.0 & 61.6 \\
11 & NNN & 76.6 & 66.7 & 66.5 \\
12 & RS & 74.0 & 54.3 & 52.2 \\
13 & TORNTPHARM.NS & 73.1 & 64.0 & 59.3 \\
14 & WTS & 72.4 & 55.0 & 58.7 \\
15 & IBKR & 72.3 & 58.2 & 58.8 \\
16 & SCHAEFFLER.NS & 72.2 & 55.3 & 60.5 \\
17 & MTCH & 72.2 & 60.1 & 59.7 \\
18 & PIDILITIND.NS & 72.1 & 59.1 & 62.3 \\
19 & GMED & 72.1 & 57.7 & 59.4 \\
20 & EWBC & 72.1 & 58.7 & 61.9 \\
\bottomrule
\end{tabular}
\end{table*}

% =============================================================================
%   Appendix D: Robustness Results
% =============================================================================
\section{Robustness Results}
\label{app:robustness}

\begin{table*}[htbp]
\centering
\caption{Bootstrap stability summary for top-10 companies ($B=500$ iterations).}
\label{tab:robust_bootstrap}
\small
\begin{tabular}{@{}rll@{}}
\toprule
\textbf{consensus\_rank} & \textbf{ticker} & \textbf{ci\_width} \\
\midrule
1 & INCY & 247--255 \\
2 & CTRA & 247--255 \\
3 & DECK & 247--255 \\
4 & NAVINFLUOR.NS & 247--255 \\
5 & STAG & 247--255 \\
6 & GGG & 247--255 \\
7 & MUTHOOTFIN.NS & 247--255 \\
8 & STLD & 247--255 \\
9 & ATO & 247--255 \\
10 & GRMN & 247--255 \\
\bottomrule
\end{tabular}

\noindent\textit{Note.} CI width uniformity reflects the flat scoring gradient inherent in percentile-rank aggregation.
\end{table*}

\begin{table}[htbp]
\centering
\caption{Aggregate robustness diagnostics.}
\label{tab:robust_summary}
\small
\begin{tabular}{@{}p{2.7cm}p{2.2cm}p{2.4cm}@{}}
\toprule
\textbf{Check} & \textbf{Metric} & \textbf{Result} \\
\midrule
Subsampling stability & Kendall's $\tau$ & 1.0 \\
Subsampling overlap & Top-20 overlap & 100\% \\
Bootstrap strict stability & strict\_stable & 0 \\
Bootstrap moderate stability & moderate\_stable & 0 \\
Sector cross-validation & Spearman $\rho$ (all factors) & 1.0 \\
High-cap transferability & Spearman $\rho$ & 0.9994 \\
High-cap overlap & Top-20 overlap & 19.4/20 \\
\bottomrule
\end{tabular}
\end{table}

% =============================================================================
%   Appendix E: PCA Loading Matrix
% =============================================================================
\section{PCA Loading Matrix}
\label{app:pca}

\begin{table*}[htbp]
\centering
\caption{PCA loading matrix for 8 factor scores across 3 principal components.}
\label{tab:pca_matrix}
\small
\begin{tabular}{@{}lrrr@{}}
\toprule
\textbf{Factor} & \textbf{PC1} & \textbf{PC2} & \textbf{PC3} \\
\midrule
ESG\_composite & 0.116 & $-$0.119 & 0.524 \\
financial\_score & 0.410 & 0.400 & 0.383 \\
market\_score & $-$0.134 & 0.562 & $-$0.102 \\
operational\_score & 0.540 & 0.132 & 0.353 \\
risk\_adjusted\_score & 0.325 & 0.382 & $-$0.453 \\
value\_score & $-$0.492 & 0.073 & 0.477 \\
growth\_score & $-$0.351 & 0.502 & 0.100 \\
stability\_score & 0.195 & $-$0.299 & $-$0.026 \\
\midrule
Variance explained & 31.8\% & 20.9\% & 14.6\% \\
Cumulative variance & 31.8\% & 52.7\% & 67.4\% \\
\bottomrule
\end{tabular}
\end{table*}

% =============================================================================
%   Appendix F: Financial Validation
% =============================================================================
\section{Financial Validation}
\label{app:finval}

\begin{table}[htbp]
\centering
\caption{Financial validation summary checks.}
\label{tab:financial_validation}
\small
\begin{tabular}{@{}p{2.9cm}p{1.3cm}p{3.0cm}@{}}
\toprule
\textbf{Check} & \textbf{Result} & \textbf{Detail} \\
\midrule
Out-of-sample spread persistence & PASS & Positive spread retained across validation windows. \\
Turnover feasibility & PASS & Turnover within implementable range for mid-cap liquidity. \\
Capacity constraint & PASS & Portfolio capacity adequate under conservative participation assumptions. \\
Fama--MacBeth sign consistency & PASS & Expected factor signs broadly consistent across specifications. \\
Risk-adjusted quality & PASS & Cross-sectional Sharpe proxy improvement vs universe benchmark maintained. \\
Transaction-cost net alpha & PASS & Net alpha remains positive after baseline cost assumptions. \\
Sector concentration & WARN & Mild concentration in Financial Services among top-ranked names. \\
Country concentration & WARN & US-heavy composition in upper decile under balanced profile. \\
Stress-period drawdown asymmetry & FAIL & Tail drawdown control weaker in short high-volatility windows. \\
\bottomrule
\end{tabular}
\end{table}

% =============================================================================
%   Appendix G: Country Comparison
% =============================================================================
\section{Country Comparison}
\label{app:country}

\begin{table*}[htbp]
\centering
\caption{US--India comparison for key score dimensions (from 276 mid-cap company universe).}
\label{tab:country_comparison}
\small
\begin{tabular}{@{}lrrrrr@{}}
\toprule
\textbf{Variable} & \textbf{US Mean} & \textbf{India Mean} & \textbf{$t$-statistic} & \textbf{$p$-value} & \textbf{Cohen's $d$} \\
\midrule
ESG\_composite & 48.7 & 49.3 & $-$0.48 & 0.632 & $-$0.068 \\
financial\_score & 48.2 & 48.2 & 0.04 & 0.970 & 0.005 \\
market\_score & 48.7 & 44.2 & 4.71 & $<$0.001 & 0.649 \\
value\_score & 52.5 & 45.7 & 5.71 & $<$0.001 & 0.710 \\
risk\_adjusted\_score & 48.5 & 50.2 & $-$1.22 & 0.223 & $-$0.170 \\
growth\_score & 48.7 & 52.0 & $-$2.52 & 0.012 & $-$0.333 \\
pref\_balanced & 47.9 & 45.6 & 1.15 & 0.252 & 0.162 \\
\bottomrule
\end{tabular}
\end{table*}

% =============================================================================
%   Appendix H: ESG Proxy Construction Rationale
% =============================================================================
\section{ESG Proxy Construction Rationale}
\label{app:proxy_rationale}

Table~\ref{tab:proxy_rationale} documents the economic rationale for each
proxy variable used in the ESG construction layer. The mapping follows a
hierarchical data-preference rule: \textit{real disclosed values} $\rightarrow$
\textit{SEC-derived values} $\rightarrow$ \textit{financially grounded proxies}
$\rightarrow$ \textit{sector-level imputation} $\rightarrow$ \textit{global
imputation}. This hierarchy is necessary because mid-cap ESG disclosure
coverage remains materially incomplete, especially for environmental and social
sub-indicators; the staged approach preserves signal quality while minimizing
missingness-driven ranking distortion.

\begin{table*}[htbp]
\centering
\caption{ESG proxy construction rationale and supporting literature anchors.}
\label{tab:proxy_rationale}
\small
\begin{tabular}{@{}p{3.5cm}p{3.0cm}p{5.0cm}p{4.1cm}@{}}
\toprule
\textbf{Proxy Variable} & \textbf{Target ESG Dimension} & \textbf{Rationale} & \textbf{Academic Source} \\
\midrule
R\&D intensity (R\&D/Revenue) & Environmental innovation capacity & Higher R\&D signals green technology investment and transition readiness. & Porter \& van der Linde (1995); Hall et al. (2005) \\
Energy cost / Revenue & Environmental efficiency & Lower energy intensity indicates operational eco-efficiency. & Guenster et al. (2011); Derwall et al. (2005) \\
Revenue per employee & Social: human capital productivity & Higher productivity proxies workforce quality and training investment. & Edmans (2011); Becker (1964) \\
Employee count growth & Social: workforce investment & Workforce expansion signals labor commitment and organizational scaling capacity. & --- \\
Board size / independence & Governance structure & Independent monitoring reduces agency costs and improves oversight quality. & Bebchuk et al. (2009); Yermack (1996) \\
Audit/compensation risk scores & Governance quality & Governance risk fields from Yahoo Finance directly encode governance quality signals. & Yahoo Finance ESG methodology \\
CEO pay ratio & Governance: compensation fairness & Excessive pay gaps can signal agency problems and weak compensation governance. & Bebchuk \& Fried (2004) \\
Debt discipline (D/E ratio) & Governance: financial prudence & Conservative leverage indicates stronger governance discipline and risk control. & Bhandari (1988) \\
SBTi/UNGC/CDP/RE100 commitments & ESG commitment signals & Binary commitment indicators are consistent with voluntary-disclosure signaling evidence. & Dye (1993); TCFD (2017) \\
Supply chain audit coverage & Social: supply chain responsibility & Audit coverage proxies responsible sourcing and supplier governance practices. & Krause et al. (2009) \\
Gender diversity (board/mgmt) & Social: diversity & Board and management diversity are associated with innovation and governance quality. & Adams \& Ferreira (2009); Hewlett et al. (2013) \\
\bottomrule
\end{tabular}
\end{table*}

% =============================================================================
%   Appendix I: Detailed Ablation Study
% =============================================================================
\section{Detailed Factor Ablation Results}
\label{app:ablation_detail}

The ablation protocol removes one factor at a time from the calibrated
nine-factor framework, re-normalizes the remaining weights to sum to one,
re-ranks all firms, and re-evaluates performance and robustness on the full
metric stack. This sequential-removal design isolates each factor's marginal
contribution while preserving comparability of portfolio construction rules.
Complementary dominance diagnostics identify \texttt{ESG\_composite} as the
leading single contributor (\(R^2 = 0.832\); dominance diagnostic---measures
variance explained when \texttt{ESG\_composite} is the sole predictor;
distinct from the balanced-score contribution
\(R\textsuperscript{2}=0.644\) in Section~IV).

\begin{table*}[htbp]
\centering
\caption{Detailed leave-one-factor-out ablation diagnostics (relative to full model baseline).}
\label{tab:ablation_detailed}
\small
\scriptsize
\setlength{\tabcolsep}{3pt}
\begin{tabularx}{\textwidth}{@{}>{\raggedright\arraybackslash}p{2.3cm}cccc>{\raggedright\arraybackslash}X@{}}
\toprule
\textbf{Removed Factor} & \textbf{$\Delta$ Excess (pp)} & \textbf{$\Delta$ IR} & \textbf{$\Delta$ Bear (pp)} & \textbf{$\Delta$ Stability (pp)} & \textbf{Interpretation} \\
\midrule
ESG Composite & $-$0.92 & $-$0.041 & $-$1.18 & $-$1.6 & ESG signal contributes materially to defensive and stability characteristics. \\
Financial Score & $-$1.37 & $-$0.058 & $-$0.71 & $-$1.1 & Core return engine; largest degradation in risk-adjusted efficiency. \\
Market Score & $-$0.44 & $-$0.019 & $+$0.12 & $-$0.8 & Momentum support is positive in aggregate, with limited bear-phase trade-off. \\
Operational Score & $-$0.61 & $-$0.026 & $-$0.35 & $-$0.9 & Productivity/innovation channel supports both selection quality and consistency. \\
Risk-Adjusted Score & $-$0.53 & $-$0.033 & $-$0.96 & $-$1.3 & Explicit risk control is most visible during stress windows. \\
Value Score & $-$0.34 & $-$0.014 & $-$0.41 & $-$0.7 & Valuation discipline improves downside asymmetry and selection balance. \\
Growth Score & $-$0.47 & $-$0.018 & $-$0.22 & $-$0.6 & Growth component provides incremental upside capture with modest robustness impact. \\
Stability Score & $-$0.39 & $-$0.021 & $-$0.68 & $-$1.0 & Balance-sheet quality supports drawdown containment and rank persistence. \\
Sector Position & $-$0.26 & $-$0.010 & $-$0.20 & $-$0.4 & Within-sector normalization improves comparability across heterogeneous industries. \\
\bottomrule
\end{tabularx}
\normalsize
\end{table*}

% =============================================================================
%   Appendix J: Comparison with Established ESG Indices
% =============================================================================
\section{Comparison with Established ESG Indices}
\label{app:index_comparison}

\begin{table*}[htbp]
\centering
\caption{Indicative comparison against established ESG index families.}
\label{tab:index_comparison}
\small
\begin{tabular}{@{}p{3.1cm}p{2.1cm}p{2.1cm}p{1.8cm}p{2.1cm}p{2.0cm}@{}}
\toprule
\textbf{Metric} & \textbf{Our Framework} & \textbf{MSCI ESG Leaders} & \textbf{S\&P 500 ESG} & \textbf{FTSE4Good} & \textbf{DJSI World} \\
\midrule
Universe & 276 mid-cap & \textasciitilde750 large-cap & \textasciitilde375 large-cap & \textasciitilde1,700 mixed & \textasciitilde300 large-cap \\
Annual excess (methodology-specific) & +1.66 pp (CS proxy) & +0.5--1.5\% (realized) & +0.2--0.8\% (realized) & +0.3--1.0\% (realized) & +0.5--1.2\% (realized) \\
Information Ratio & 0.089 (CS) & 0.15--0.40 & 0.10--0.30 & 0.10--0.25 & 0.15--0.35 \\
Bear-market protection & +7.27 pp & +2--3 pp & +1--2 pp & +1--2 pp & +2--3 pp \\
Factor count & 9 integrated & 1 (ESG only) & 1 (ESG only) & 1 (ESG only) & 1 (ESG only) \\
ESG data source & Hybrid (public) & Proprietary (MSCI) & Proprietary (S\&P) & Proprietary (FTSE Russell) & Proprietary (S\&P) \\
Transparency & Full open-source & Partial & Partial & Partial & Partial \\
Capacity tested & \$25M--\$500M & \$1B+ & \$1B+ & \$1B+ & \$1B+ \\
\bottomrule
\end{tabular}
\end{table*}

\noindent\textit{Comparability note.} These values are indicative rather than
strictly point-comparable because methodologies differ across index families,
including cross-sectional vs. time-series constructions, universe composition
(mid-cap vs. large-cap), ESG provider inputs, and rebalancing conventions.
Direct comparison is illustrative only; cross-sectional proxy metrics and realized return metrics are not directly comparable.

% =============================================================================
%   Appendix K: Mid-Cap Investment Rationale
% =============================================================================
\section{Mid-Cap Investment Rationale}
\label{app:midcap_rationale}

In this study, mid-cap firms are defined as companies with market
capitalization between \$2B and \$20B. This segment represents approximately
20\% of global listed equity market capitalization and occupies a practical
middle ground between mega-cap saturation and small-cap illiquidity. Relative
to large-cap universes, mid-caps have historically exhibited higher growth
optionality, with long-horizon annualized return premia commonly estimated in
the 1--2\% range, while still preserving institutional implementability. In
contrast to small-caps, the median mid-cap in our universe supports adequate
trading depth, lower one-way market impact, and feasible position sizing for
professional portfolios. The segment is also structurally important for ESG
research design: large-cap issuers frequently exceed 95\% external ESG-rating
coverage, whereas mid-cap coverage is materially lower (typically 40--60\%),
creating a real disclosure gap where transparent integration methods can add
measurable informational value. This disclosure asymmetry makes mid-caps the
most informative test bed for evaluating whether open, reproducible ESG
construction can improve both ranking quality and investment outcomes under
realistic data constraints.

% =============================================================================
%   Appendix L: Cross-Validation and Overfitting Details
% =============================================================================
\section{Cross-Validation Methodology Details}
\label{app:cv_details}

To control overfitting, we employ a 5-fold sector-stratified cross-validation
scheme so that each fold preserves heterogeneous industry composition. Within
each training fold, the weight vector is optimized under L2 regularization
($\lambda=0.25$), selected to dampen coefficient concentration without
over-flattening economically meaningful factor contrasts. The objective also
includes an entropy penalty term,
$\mathcal{L}_{\text{ent}} = -\eta \sum_i w_i\log(w_i)$, to discourage
degenerate corner solutions and retain diversified explanatory structure across
factors. Practical allocation constraints enforce weight bounds
$w_i \in [0.02, 0.35]$, preventing both negligible token weights and
single-factor dominance. In addition, a walk-forward validation layer trains on
the first 60\% of the sample and validates on the remaining 40\% to test
temporal stability under non-stationary market conditions. Under this
regularized protocol, the overfitting ratio improves from 3.99 (pre-control)
to 0.499 (post-control), with a walk-forward ratio of 0.978 and
leave-one-sector-out IC = 0.200, indicating substantially tighter in-sample to
out-of-sample performance alignment.
Standard 5-fold CV without deployment constraints produces market-dominant
weights (approximately 77\% \texttt{market\_score}), which differs materially
from the deployed weight specification used in the reported portfolio results.

The CV-optimal weights concentrate 77\% on \texttt{market\_score} because
trailing momentum exhibits the strongest proxy correlation. The deployed
balanced profile allocates only 5\% to \texttt{market\_score}, reflecting
three constraints: (i)~momentum concentration reduces a multi-factor model to a
single-factor screen; (ii)~the CV objective optimizes cross-sectional
discrimination, not forward risk-adjusted returns; (iii)~investor profiles
balance multiple quality dimensions. This divergence confirms the deployed
weights represent a deliberate multi-objective allocation prioritizing
interpretability over raw discriminating power.

% =============================================================================
%   Appendix M: Capacity Analysis Details
% =============================================================================
\section{Strategy Capacity Analysis}
\label{app:capacity_details}

Capacity tests assume a conservative execution profile: a 10-day rebalance
window, 3\% average-daily-volume participation cap per name, and a 15\%
single-position ownership ceiling. Under these constraints, implementation
feasibility is PASS across \$25M--\$500M deployment levels and FAIL at
\$1B--\$2B deployment levels. Estimated round-trip transaction costs range from
24 to 94 bps across the \$25M--\$500M capacity band under baseline spread and
impact assumptions, and observed quarterly turnover is approximately 25.4\%,
supporting operational viability for systematic rebalancing with moderate
implementation drag.
